% The whole SMC-DSDP blueprint: Foundations, Corrupted Alice, Corrupted relay.
% web.tex inputs this file and nothing else.
%
% Each security chapter states one headline theorem and chains it to its leaves:
% a standing-hypotheses block (what the bound assumes), the statement, its direct
% dependencies as nodes, and the library/plumbing facts it rests on cited as
% black-box leaves in prose. Node bodies are terse declarative statements, per
% the mathcomp-qbs standard.

\chapter{The corrupted-Alice and corrupted-relay security of DSDP}

DSDP is a three-party protocol that computes the scalar product of private input
vectors held by Alice, Bob, and Charlie without revealing them. This blueprint
presents its security against a single semi-honest corrupted party, proved in
Rocq. The analysis has three threat models: a corrupted Alice, who both receives
the ciphertext traffic and decrypts the output; a corrupted relay (Bob or
Charlie), who forwards masked ciphertexts; and, as the tightness boundary of the
corrupted-Alice case, a corrupted Alice who fixes her query to a coordinate
vector and reads a relay's input directly.

Every headline theorem is proved in one of the axis files under
\texttt{dsdp/counting/} and \texttt{dsdp/hopping/}, over the scheme-independent
game vocabulary of \texttt{computational\_security/}, and read at the three
instance sequences of \texttt{dsdp\_instance\_sequence.v} and at the uniform
counting witness of \texttt{dsdp\_random\_inputs.v}. Each chapter below takes
one of them and chains it
to its leaves, opening with the hypotheses the bound actually assumes so an
over-hypothesised or vacuous bound is visible on the page. The two proof
techniques are kept distinct: the corrupted-Alice ciphertext channel is bounded
by a two-hop ladder over one explicit Infotheo distribution, and the output and
relay channels are bounded information-theoretically by conditional entropy and
fiber-counting.

Every $\varepsilon$ in the hop ladder is the real-or-zero advantage of a
reduction constructed in the same file as the bound that names it, so the
computational leg carries no cryptographic assumption as an axiom. Chapter
\ref{ch:adversary_class} adds the one assumed quantity, the advantage an
adversary-class assumption assigns to its class, and shows what forces such an
assumption to restrict its class. At the two concrete schemes that advantage is
itself derived, from the $e$-th residuosity assumption, so what stays assumed
at Paillier and at Benaloh is residuosity.

\part{Foundations}

\chapter{The protocol, its views, and the threat models}

The protocol runs on the ring $\mathbb{Z}_{pq}$ with $p,q$ distinct primes, so the
message space $\mathsf{msg} = \mathbb{Z}_{pq}$ has $m = pq$ elements. Each party
holds a private input and a private multiplier; the protocol computes
$S = u_1 v_1 + u_2 v_2 + u_3 v_3$ and returns it to Alice.

\begin{definition}[The message flow]
  \label{def:protocol_flow}
  \rocqok
  % library: the party programs palice/pbob/pcharlie live in core/dsdp_pismc.v,
  %   cited in prose (not in the coqdoc MODULES scope).
  The three parties exchange messages in a ring. Bob and Charlie each send Alice a
  head ciphertext of their own input. Alice sends Bob two homomorphic assemblies
  $a_1, a_2$; Bob decrypts, combines them with Charlie's contribution, and
  forwards the aggregate to Charlie; Charlie decrypts, re-encrypts the total under
  Alice's key, and returns it. Alice decrypts and removes her masks to read $S$.
  \[ \text{Alice} \to \text{Bob} \to \text{Charlie} \to \text{Alice} \]
\end{definition}

\begin{definition}[Party views]
  \label{def:party_views}
  \rocqok
  A corrupted party's view is everything it sees: its own key and inputs, the
  ciphertexts it receives, and, for Alice, the decrypted output $S$. Alice's view
  carries her key, her inputs and masks, the output $S$, and the three ciphertext
  hops; a relay's view carries its key, its own input, and the ciphertexts it
  forwards.
\end{definition}

\begin{lemma}[Protocol correctness]
  \label{lem:dsdp_is_correct}
  \rocqok
  % library: dsdp_is_correct (core/dsdp_correctness.v), dsdp_result_correct
  %   (counting/dsdp_entropy_trace.v). Cited in prose; not in MODULES scope.
  The protocol computes the scalar product: the value Alice reads equals
  $u_1 v_1 + u_2 v_2 + u_3 v_3$. Security without this is empty, since a protocol
  returning a constant is trivially private; correctness is what makes the private
  quantity the intended one.
  \[ \text{Alice's output} = u_1 v_1 + u_2 v_2 + u_3 v_3 \]
\end{lemma}

\part{Corrupted Alice}

Alice is the party the analysis works hardest on: she receives every ciphertext
and she decrypts the output. Her secrecy has two channels. The output she decrypts
is the scalar product, and conditioned on it a relay's input stays uniform over a
fiber of $m$ points, an information-theoretic bound. The ciphertexts she receives
are bounded by a ladder of two hops, each replacing one received ciphertext by an
encryption of zero. The two channels compose to the bound
$1/m + \epshopsum{D_g}$ on any predictor $g$ of Bob's input, where each summand is
the real-or-zero advantage, at the key of the party that sent that ciphertext, of
the reduction $g$ induces on that hop. The part closes with the tightness
boundary: the fiber bound needs Alice's third weight to act invertibly, and the
degenerate query that violates it lets Alice read a relay's input outright.

\chapter{What the output alone reveals: the plaintext residual}

\begin{theorem}[Residual uncertainty of the relay inputs]
  \label{thm:dsdp_centropy_uniform}
  \rocq{infotheo.dumas2017dual.dsdp.counting.dsdp_entropy.dsdp_centropy_uniform,
       infotheo.dumas2017dual.dsdp.counting.dsdp_random_inputs.dsdp_random_inputs}
  \rocqok
  \uses{def:party_views,lem:dsdp_is_correct}
  Conditioned on Alice's inputs and the decrypted output $(V_1,U_1,U_2,U_3,S)$ ---
  the plaintext residual, not her full ciphertext-carrying view --- and assuming
  the relay pair $(V_2,V_3)$ is a priori uniform and independent of the inputs, the
  relay inputs retain $\log m$ bits of uncertainty, provided Alice's third weight
  satisfies $0 < U_3 < \min(p,q)$.
  \[ \HH(V_2, V_3 \mid V_1,U_1,U_2,U_3,S) = \log m \]
  Proved in \texttt{dsdp\_entropy.v} over the two primes, the sample law and
  the eleven random inputs, with their uniformity and their independence
  assumed and the conditioner's $S$ computed from the inputs rather than left
  free. Those hypotheses are exactly the fields of the record
  \texttt{dsdp\_random\_inputs}, which \texttt{uniform\_inputs} inhabits at
  every modulus. Because the each-against-the-rest independence reaches
  Alice's three weights, the equality covers honest sampling, where her
  weights are independent of her input and of one another. At
  every modulus the equality is exact and holds against an adversary of any
  running time.
\end{theorem}
\begin{proof}
  \uses{lem:dsdp_fiber_uniform}
  \rocqok
  Each conditioning value fixes a fiber of the scalar-product constraint on which
  $(V_2,V_3)$ is uniform (Lemma~\ref{lem:dsdp_fiber_uniform}); the conditional
  entropy is then $\log$ of the fiber size, which is $m$.
\end{proof}

\begin{lemma}[Fiber uniformity of the DSDP solution]
  \label{lem:dsdp_fiber_uniform}
  \rocq{infotheo.dumas2017dual.dsdp.counting.dsdp_entropy.Pr_dsdp_sol_uniform_ring,
       infotheo.dumas2017dual.dsdp.counting.dsdp_entropy.dsdp_fiber_card_ring}
  \rocqok
  When the multiplier $U_3$ acts injectively, the solution set of the DSDP
  constraint has exactly $m$ points, and each solution has conditional probability
  $1/m$ given the inputs and the output. This is the counting fact both the
  entropy bound and the guessing bound consume.
  \[ \#\lvert\mathrm{fiber}\rvert = m,\qquad
        \Pr[(V_2,V_3)=(v_2,v_3)\mid \mathrm{cond}] = 1/m \]
\end{lemma}
\begin{proof}
  \rocqok
  Over a commutative ring the fiber is the image of the ring under $v \mapsto U_3
  v$, of size $\#R = m$; conditional uniformity follows from the joint uniformity
  and independence hypotheses (the route-F Infotheo fiber theorems).
\end{proof}

\begin{theorem}[N-party residual]
  \label{thm:dsdp_centropy_uniform_n}
  \rocq{infotheo.dumas2017dual.dsdp.counting.dsdp_entropy.dsdp_centropy_uniform_n}
  \rocqok
  The same argument at $n$ relays: conditioned on the N-party view, the relay
  inputs retain $\log(m^n)$ bits, one full message of uncertainty per relay
  dimension. It is a generic lemma over abstract variables satisfying the fiber,
  projection, uniformity, independence, and joint-marginal hypotheses; no N-party
  DSDP program instance is supplied.
  \[ \HH(V_0,\dots,V_{n} \mid \mathrm{view}) = \log(m^{n}) \]
\end{theorem}

\section{Reused Infotheo facts}

The fiber lemma of this chapter and the endpoint bound of the next one are
assembled from four facts that live outside the DSDP files: two counting facts
about the solution set of the scalar-product constraint, one conditional
uniformity fact for a fiber, and one diagonal bound for conditionally independent
variables. They are recorded here as leaves.

\begin{lemma}[Fiber cardinality is $m$]
  \label{lem:exist_fiber_card}
  \rocq{infotheo.dumas2017dual.dsdp.counting.dsdp_entropy.dsdp_fiber_card_ring}
  \rocqok
  Over a commutative ring, when the multiplier $u_3$ is injective the solution
  set of the scalar-product constraint has exactly $\#R = m$ points, the image of
  the ring under $v \mapsto u_3 v$.
  \[ \#\lvert \fib(u_1,u_2,u_3,v_1,s)\rvert = \#R = m \]
\end{lemma}

\begin{lemma}[Conditional uniformity on the fiber]
  \label{lem:exist_cpr_fiber}
  \rocq{infotheo.dumas2017dual.entropy_fiber.entropy_fiber_zpq.gen_cPr_uniform_fiber}
  \rocqok
  \uses{lem:exist_fiber_card}
  For a joint distribution in which the variable pair is uniform and
  independent of the inputs, conditioning on the revealed values leaves the
  variable uniform over its fiber.
  \[ \Pr[\,\mathrm{Var}=v \mid \mathrm{Cond}=\cond\,]
        = \#\lvert\fib(\cond)\rvert^{-1} \]
\end{lemma}

\begin{lemma}[DSDP solution is uniform, probability $1/m$]
  \label{lem:exist_pr_sol}
  \rocq{infotheo.dumas2017dual.dsdp.counting.dsdp_entropy.Pr_dsdp_sol_uniform_ring}
  \rocqok
  \uses{lem:exist_cpr_fiber,lem:exist_fiber_card}
  Over a commutative ring with injective multiplier, each solution $(v_2,v_3)$ of
  the DSDP constraint has conditional probability $1/m$ given the inputs and $S$.
  \[ \Pr[\,(V_2,V_3)=(v_2,v_3)\mid \mathrm{Cond}=\cond\,] = m^{-1} \]
\end{lemma}

\begin{lemma}[Conditional-independence diagonal bound]
  \label{lem:cinde_diagonal_bound}
  \rocq{infotheo.dumas2017dual.lib.extra_proba.cinde_diagonal_bound}
  \rocqok
  When $X$ and $Y$ are conditionally independent given $Z$ and each value of $Y$
  has conditional probability at most $1/m$ given $Z$, the probability of the
  diagonal $X = Y$ is at most $1/m$.
  \[ X \perp Y \mid Z,\quad \Pr[\,Y = a \mid Z = c\,] \le 1/m
        \;\Longrightarrow\; \Pr[\,X = Y\,] \le 1/m \]
\end{lemma}

\chapter{Computational secrecy of the ciphertext channel}
\label{ch:hopping}

The two ciphertexts Alice receives are bounded by a ladder of two hops over one
explicit distribution. Hop $i$ replaces the $i$-th received ciphertext by an
encryption of zero, and the loss it incurs is exactly the
real-or-zero advantage of a single-query adversary constructed in the same file.
At the end of the ladder both slots encrypt zero and the fiber bound applies.
Throughout, $m$ is the size of the plaintext space and the plaintext carrier is a
finite commutative unit ring.

\begin{definition}[The corrupted-Alice experiment]
  \label{def:alice_experiment}
  \rocq{infotheo.dumas2017dual.dsdp.hopping.dsdp_alice_hop_secrecy.alice_sampleT,
       infotheo.dumas2017dual.dsdp.hopping.dsdp_alice_hop_secrecy.alice_sample_fdist,
       infotheo.dumas2017dual.dsdp.hopping.dsdp_alice_hop_secrecy.alice_sample_fdistE,
       infotheo.dumas2017dual.dsdp.hopping.dsdp_alice_hop_secrecy.sample_V2,
       infotheo.dumas2017dual.dsdp.hopping.dsdp_alice_hop_secrecy.sample_V3,
       infotheo.dumas2017dual.dsdp.hopping.dsdp_alice_hop_secrecy.sample_R2,
       infotheo.dumas2017dual.dsdp.hopping.dsdp_alice_hop_secrecy.sample_R3,
       infotheo.dumas2017dual.dsdp.hopping.dsdp_alice_hop_secrecy.Rho2,
       infotheo.dumas2017dual.dsdp.hopping.dsdp_alice_hop_secrecy.Rho3,
       infotheo.dumas2017dual.dsdp.hopping.dsdp_alice_hop_secrecy.RA1,
       infotheo.dumas2017dual.dsdp.hopping.dsdp_alice_hop_secrecy.RA2,
       infotheo.dumas2017dual.dsdp.hopping.dsdp_alice_hop_secrecy.Sout,
       infotheo.dumas2017dual.dsdp.hopping.dsdp_alice_hop_secrecy.SoutE}
  \rocqok
  \uses{def:party_views}
  The experiment is the uniform distribution on the product of eight
  coordinates: the two honest inputs $V_2, V_3$, Alice's two mask plaintexts
  $R_2, R_3$, the randomness $\rho_2, \rho_3$ of the two ciphertexts she
  receives, and the randomness $\alpha_1, \alpha_2$ of her two outgoing combines.
  Alice's own input $v_1$ and her three weights $u_1, u_2, u_3$ are parameters of
  the experiment rather than sampled coordinates. The output she learns is the
  weighted scalar product of her weights with the two honest inputs.
  \[ S = u_1 v_1 + u_2 V_2 + u_3 V_3 \]
\end{definition}

\begin{definition}[The hop ladder]
  \label{def:alice_hop_ladder}
  \rocq{infotheo.dumas2017dual.dsdp.hopping.dsdp_alice_hop_secrecy.bob_real_cipher,
       infotheo.dumas2017dual.dsdp.hopping.dsdp_alice_hop_secrecy.bob_zero_cipher,
       infotheo.dumas2017dual.dsdp.hopping.dsdp_alice_hop_secrecy.charlie_real_cipher,
       infotheo.dumas2017dual.dsdp.hopping.dsdp_alice_hop_secrecy.charlie_zero_cipher,
       infotheo.dumas2017dual.dsdp.hopping.dsdp_alice_hop_secrecy.alice_hop_tupleT,
       infotheo.dumas2017dual.dsdp.hopping.dsdp_alice_hop_secrecy.alice_tuple_real,
       infotheo.dumas2017dual.dsdp.hopping.dsdp_alice_hop_secrecy.alice_tuple_bob_zero,
       infotheo.dumas2017dual.dsdp.hopping.dsdp_alice_hop_secrecy.alice_tuple_all_zero,
       infotheo.computational_security.indcpa_game.accept,
       infotheo.computational_security.indcpa_game.acceptE,
       infotheo.computational_security.indcpa_game.accept_ge0}
  \rocqok
  \uses{def:alice_experiment}
  Alice's tuple carries her two masks, her two combine randomnesses, the
  output $S$, and the two ciphertexts she receives. Three tuples differ only
  in the plaintexts of those two ciphertexts.
  \[ \hoptuple{0} = \big(R_2, R_3, \alpha_1, \alpha_2, S,\;
       \Enc(\pk_B, V_2, \rho_2),\; \Enc(\pk_C, V_3, \rho_3)\big) \]
  \[ \hoptuple{1} = \big(R_2, R_3, \alpha_1, \alpha_2, S,\;
       \Enc(\pk_B, 0, \rho_2),\; \Enc(\pk_C, V_3, \rho_3)\big) \]
  \[ \hoptuple{2} = \big(R_2, R_3, \alpha_1, \alpha_2, S,\;
       \Enc(\pk_B, 0, \rho_2),\; \Enc(\pk_C, 0, \rho_3)\big) \]
  $\hoptuple{0}$ is Alice's real tuple and $\hoptuple{2}$ the all-zero
  endpoint. $G_i$ is the joint law of the honest input pair beside
  $\hoptuple{i}$, and $\gamma_i(D)$ the probability that a distinguisher $D$
  accepts a value sampled from $G_i$.
\end{definition}

\begin{definition}[Real-or-zero advantage of a single-query adversary]
  \label{def:indcpa_epsilon}
  \rocq{infotheo.computational_security.indcpa_game.enc_fdist,
       infotheo.computational_security.indcpa_game.indcpa_adversary,
       infotheo.computational_security.indcpa_game.indcpa_challenger,
       infotheo.computational_security.indcpa_game.indcpa_experiment,
       infotheo.computational_security.indcpa_game.indcpa_accept,
       infotheo.computational_security.indcpa_game.indcpa_success_real,
       infotheo.computational_security.indcpa_game.indcpa_success_realE,
       infotheo.computational_security.indcpa_game.indcpa_success_zero,
       infotheo.computational_security.indcpa_game.indcpa_success_zeroE,
       infotheo.computational_security.indcpa_game.indcpa_epsilon}
  \rocqok
  A single-query real-or-zero adversary $A$ is a finite state type, a law on it,
  a challenge plaintext $p_A$ read off the state, and a decision taken on the
  state and one challenge ciphertext. Encryption randomness is uniform. The
  challenger at hidden bit $b$ is the encryption law of $p_A$ at $b = 1$ and the
  encryption law of zero at $b = 0$, and the experiment at $b$ is the law of the
  adversary's decision against that challenge. The advantage of $A$ at a public
  key is the gap between its acceptance probability at the two bits.
  \[ \epsror{\pk}{A} = \big\lvert \Pr[A \text{ accepts } \Enc(\pk, p_A)]
       - \Pr[A \text{ accepts } \Enc(\pk, 0)] \big\rvert \]
\end{definition}

\begin{definition}[The two hop reductions]
  \label{def:alice_hop_reductions}
  \rocq{infotheo.dumas2017dual.dsdp.hopping.dsdp_alice_hop_secrecy.hop0_stateT,
       infotheo.dumas2017dual.dsdp.hopping.dsdp_alice_hop_secrecy.Hop0State,
       infotheo.dumas2017dual.dsdp.hopping.dsdp_alice_hop_secrecy.hop0_assemble,
       infotheo.dumas2017dual.dsdp.hopping.dsdp_alice_main.bob_challenge_adversary,
       infotheo.dumas2017dual.dsdp.hopping.dsdp_alice_hop_secrecy.hop1_stateT,
       infotheo.dumas2017dual.dsdp.hopping.dsdp_alice_hop_secrecy.Hop1State,
       infotheo.dumas2017dual.dsdp.hopping.dsdp_alice_hop_secrecy.hop1_assemble,
       infotheo.dumas2017dual.dsdp.hopping.dsdp_alice_main.charlie_challenge_adversary}
  \rocqok
  \uses{def:alice_hop_ladder,def:indcpa_epsilon}
  The hop-$i$ reduction $\hopred{i}{D}$ of a distinguisher $D$ is the
  single-query adversary whose state carries everything the hop-$i$ tuple holds
  apart from the $i$-th ciphertext slot, whose challenge plaintext is the honest
  input that slot carries, and whose decision is $D$ applied to the honest input
  pair together with the tuple rebuilt around the challenge. Hop $0$ challenges
  Bob's key on $V_2$, and hop $1$ challenges Charlie's key on $V_3$.
\end{definition}

\paragraph{How a hop reduction is used.}
A hop reduction converts a distinguisher for two adjacent protocol hops into
a real-or-zero adversary for the ciphertext slot changed by that hop. For the
first hop, the reduction retains every other component of the sampled joint
value, places Bob's challenge ciphertext in the missing slot, and applies
$D$ to the rebuilt value. The construction can be summarized as follows.
\[
\begin{array}{c}
D \text{ distinguishing hop $0$ from hop $1$}
\\[0.6ex]
\Big\downarrow\ {\scriptstyle D \mapsto \hopred{0}{D}}
\\[0.6ex]
\hopred{0}{D} \text{ distinguishing }
\Enc(\pk_B,V_2) \text{ from } \Enc(\pk_B,0)
\end{array}
\]
The bottom line abbreviates the complete real-or-zero experiment. The
adversary $\hopred{0}{D}$ receives its sampled state together with the
challenge ciphertext and reconstructs the full input expected by $D$. When
the challenge encrypts $V_2$, that input has the law of hop $0$. When it
encrypts zero, the input has the law of hop $1$. Consequently, the two
distinguishing gaps are equal. The second reduction applies the same
construction to hops $1$ and $2$, using Charlie's key and challenge plaintext
$V_3$.

\begin{lemma}[One hop loses one reduction advantage]
  \label{lem:alice_view_advantage}
  \rocq{infotheo.dumas2017dual.dsdp.hopping.dsdp_alice_main.hop0_advantageE,
       infotheo.dumas2017dual.dsdp.hopping.dsdp_alice_main.hop1_advantageE}
  \rocqok
  \uses{def:alice_hop_reductions}
  For all weights $u_1, u_2, u_3$ and every distinguisher $D$, zeroing one
  received ciphertext moves the acceptance probability by exactly the advantage
  of the corresponding hop reduction, at the key of the party that sent that
  ciphertext.
  \[ \lvert \gamma_0(D) - \gamma_1(D)\rvert = \epsror{\pk_B}{\hopred{0}{D}},
     \qquad
     \lvert \gamma_1(D) - \gamma_2(D)\rvert = \epsror{\pk_C}{\hopred{1}{D}} \]
\end{lemma}

\begin{lemma}[Guessing at the all-zero endpoint]
  \label{lem:guess_all_zero}
  \rocq{infotheo.dumas2017dual.dsdp.hopping.dsdp_alice_hop_secrecy.AliceSpectator,
       infotheo.dumas2017dual.dsdp.hopping.dsdp_alice_hop_secrecy.alice_spectator_of,
       infotheo.dumas2017dual.dsdp.hopping.dsdp_alice_hop_secrecy.alice_spectator_indep,
       infotheo.dumas2017dual.dsdp.hopping.dsdp_alice_hop_secrecy.alice_spectator_cinde,
       infotheo.dumas2017dual.dsdp.hopping.dsdp_alice_hop_secrecy.alice_VarRV_cond_uniform,
       infotheo.dumas2017dual.dsdp.hopping.dsdp_alice_hop_secrecy.alice_V2_cond_Sout,
       infotheo.dumas2017dual.dsdp.hopping.dsdp_alice_hop_secrecy.alice_V2_cond_le,
       infotheo.dumas2017dual.dsdp.hopping.dsdp_alice_hop_secrecy.alice_hop_tuple_of_spectator,
       infotheo.dumas2017dual.dsdp.hopping.dsdp_alice_hop_secrecy.all_zero_guess_V2_le_invm}
  \rocqok
  \uses{def:alice_hop_ladder,lem:exist_pr_sol,lem:cinde_diagonal_bound}
  When Alice's third weight $u_3$ is a unit of the plaintext ring, every
  predictor reading the all-zero tuple matches Bob's input with probability at
  most $1/m$. The all-zero tuple is independent of the honest inputs, hence
  conditionally independent of $V_2$ given the output, and conditioned on the
  output alone $V_2$ is uniform.
  \[ u_3 \in R^{\times} \;\Longrightarrow\;
       \Pr[\,g(\hoptuple{2}) = V_2\,] \le 1/m \]
\end{lemma}

\begin{definition}[The distinguisher of a predictor]
  \label{def:distinguisher_of_predictor}
  \rocq{infotheo.computational_security.indcpa_game.distinguisher_of_predictor,
       infotheo.dumas2017dual.dsdp.hopping.dsdp_alice_main.guess_V2_acceptE}
  \rocqok
  \uses{def:alice_hop_ladder}
  The distinguisher $D_g$ of a predictor $g$ accepts when $g$, applied to the
  tuple component of its input, returns the first honest input. The event that
  $g$ matches Bob's input at hop $i$ is the acceptance event of $D_g$ at hop $i$.
  \[ D_g(v_2, v_3, t) = [\,g(t) = v_2\,] \]
\end{definition}

\begin{theorem}[Guessing Bob's input from Alice's real tuple]
  \label{thm:alice_guess_real}
  \rocq{infotheo.dumas2017dual.dsdp.hopping.dsdp_alice_main.alice_tuple_guess_V2_le}
  \rocqok
  \uses{thm:alice_simulation_secure,lem:simulator_factorization,lem:guess_all_zero,def:distinguisher_of_predictor}
  When $u_3$ is a unit of the plaintext ring, every predictor reading Alice's
  real tuple matches Bob's input with probability at most $1/m$ plus the
  advantages of the two hop reductions of its distinguisher.
  \[ \Pr[\,g(\hoptuple{0}) = V_2\,]
       \le \tfrac{1}{m} + \epshopsum{D_g} \]
  The bound is the simulation bound of
  Theorem~\ref{thm:alice_simulation_secure} at the predictor's own
  distinguisher $D_g$, which accepts exactly when $g$ returns Bob's input, plus
  what the simulator's output reveals: on the ideal-world law $g$ succeeds
  with probability at most $1/m$, the mass of the fiber of the leaked output
  (Lemma~\ref{lem:guess_all_zero}), and the real law is within the two hop
  advantages of it.
  Proved in \texttt{dsdp\_alice\_main.v} as
  \texttt{alice\_tuple\_guess\_V2\_le}, over the data of one DSDP scheme, so a
  reading at a concrete scheme supplies that data and proves nothing further.
\end{theorem}
\begin{proof}
  \uses{thm:alice_simulation_secure,lem:simulator_factorization,lem:guess_all_zero}
  \rocqok
  $\Pr_{\mathrm{real}}[D_g] \le \Pr_{\mathrm{ideal}}[D_g]
   + |\Pr_{\mathrm{real}}[D_g] - \Pr_{\mathrm{ideal}}[D_g]|$; the first term
  is the all-zero endpoint probability, at most $1/m$, and the second is the
  simulation bound.
\end{proof}

\chapter{Simulation-based security}
\label{ch:simulation}

Alice's tuple can be produced by a simulator that reads the output and nothing
else, up to the same two-reduction sum. The scope is average-case: the honest
inputs are sampled uniformly by the experiment on one side and fed to the ideal
functionality on the other, rather than quantified over all input vectors.

\begin{definition}[The simulator and the ideal-world law]
  \label{def:alice_simulator}
  \rocq{infotheo.dumas2017dual.dsdp.hopping.dsdp_alice_hop_secrecy.alice_simulator,
       infotheo.dumas2017dual.dsdp.hopping.dsdp_alice_hop_secrecy.alice_ideal}
  \rocqok
  \uses{def:alice_hop_ladder}
  From a value $s$ of the output, the simulator draws uniform masks, uniform
  combine randomness, and one encryption of zero under each of Bob's and
  Charlie's keys, and returns $s$ unchanged. The ideal-world law samples the
  honest input pair, computes its output, and returns that pair with a simulated
  tuple.
  \[ \Sim(s) = \big(\mathcal{U}, \mathcal{U}, s,
       \Enc(\pk_B, 0), \Enc(\pk_C, 0)\big),\qquad
     \Ideal = \big[(v_2,v_3) \sim \mathcal{U};\ t \sim \Sim(S(v_2,v_3));\
       (v_2, v_3, t)\big] \]
\end{definition}

\begin{lemma}[The all-zero endpoint is the simulator]
  \label{lem:simulator_factorization}
  \rocq{infotheo.dumas2017dual.dsdp.hopping.dsdp_alice_hop_secrecy.alice_spectator_law,
       infotheo.dumas2017dual.dsdp.hopping.dsdp_alice_hop_secrecy.alice_hop_tuple_all_zero_pfwd1E,
       infotheo.dumas2017dual.dsdp.hopping.dsdp_alice_hop_secrecy.dsdp_alice_hop_tuple_cond_sim,
       infotheo.dumas2017dual.dsdp.hopping.dsdp_alice_main.alice_idealE}
  \rocqok
  \uses{def:alice_simulator}
  Conditioned on any honest input pair of positive probability, Alice's all-zero
  tuple has exactly the simulator law fed the output of that pair. The
  ideal-world law is therefore the joint law of the honest inputs and the
  all-zero tuple, with no error term.
  \[ \Pr[\,\hoptuple{2} = t \mid (V_2,V_3) = (v_2,v_3)\,] = \Sim(S(v_2,v_3))(t),
     \qquad \Ideal = \mathrm{law}(V_2, V_3, \hoptuple{2}) \]
\end{lemma}

\begin{theorem}[Bounded simulation security]
  \label{thm:alice_simulation_secure}
  \rocq{infotheo.dumas2017dual.dsdp.hopping.dsdp_alice_main.alice_sim_advantage_le}
  \rocqok
  \uses{def:party_views,lem:alice_view_advantage,lem:simulator_factorization}
  For all weights $u_1, u_2, u_3$, every distinguisher separates the real joint
  law of the honest inputs and Alice's tuple from the ideal-world law by at most
  the sum of the advantages of its two hop reductions.
  \[ \big\lvert \Pr_{\mathrm{law}(V_2,V_3,\hoptuple{0})}[D]
       - \Pr_{\Ideal}[D] \big\rvert \le \epshopsum{D} \]
  Proved in \texttt{dsdp\_alice\_main.v} as
  \texttt{alice\_sim\_advantage\_le}, over the data of one DSDP scheme.
\end{theorem}
\begin{proof}
  \uses{lem:simulator_factorization,lem:alice_view_advantage}
  \rocqok
  The ideal-world law is the hop-$2$ law, so the gap is
  $\lvert \gamma_0(D) - \gamma_2(D) \rvert$. The triangle inequality splits it
  at $\gamma_1(D)$, and Lemma~\ref{lem:alice_view_advantage} reads each of the
  two halves as the advantage of the reduction at one key.
\end{proof}

\begin{remark}[Worst-case simulation: open]
  \label{rem:worstcase_sim}
  \rocqok
  % blue frontier node: the standard, input-quantified simulation statement is not
  %   proved at hop 0; only the uniform-average is. Recorded to make the gap visible.
  The standard simulation-security statement quantifies over every input pair:
  for all $v_2, v_3$, the real tuple is indistinguishable from the simulated
  ideal. Theorem~\ref{thm:alice_simulation_secure} bounds the gap at the joint
  law in which the honest inputs are sampled uniformly, and the per-input
  statement is proved at the all-zero endpoint alone
  (Lemma~\ref{lem:simulator_factorization}). At hop $0$ it is open.
\end{remark}

\chapter{The bounds at the executed protocol trace}
\label{ch:trace_link}

The hopping tuple is an abstraction of what Alice sees. Identifying it with the
trace the piSMC interpreter returns when it runs the three DSDP party programs
carries the three bounds above to the executed protocol.

\begin{definition}[Alice's executed trace]
  \label{def:alice_trace}
  \rocq{infotheo.dumas2017dual.dsdp.hopping.dsdp_alice_trace_link.trace_dataT,
       infotheo.dumas2017dual.dsdp.hopping.dsdp_alice_trace_link.trace_data_of_di_data,
       infotheo.dumas2017dual.dsdp.hopping.dsdp_alice_trace_link.dsdp_procs_std,
       infotheo.dumas2017dual.dsdp.hopping.dsdp_alice_trace_link.dsdp_protocol,
       infotheo.dumas2017dual.dsdp.hopping.dsdp_alice_trace_link.trace_of_run,
       infotheo.dumas2017dual.dsdp.hopping.dsdp_alice_trace_link.AliceTrace,
       infotheo.dumas2017dual.dsdp.hopping.dsdp_alice_trace_link.alice_trace_of_hop_tuple}
  \rocqok
  \uses{def:protocol_flow,def:alice_hop_ladder}
  The DSDP protocol is the triple of party programs at the coordinates of one
  sample of the experiment, and $\mathrm{trace\_of\_run}(\mathit{ps}, i)$ is the
  message log the interpreter returns for party $i$ when it runs the process
  list $\mathit{ps}$ at a sample, encoded as a bounded sequence of at most
  fifteen entries, each a plaintext, a ciphertext, or a private key. Alice's
  trace is that reader at the DSDP protocol and at Alice.
  \[ \mathrm{AliceTrace} = \mathrm{trace\_of\_run}(\mathrm{dsdp\_protocol},
       \mathrm{Alice}) \]
\end{definition}

\begin{lemma}[The trace is a function of the tuple]
  \label{lem:trace_of_hop_tuple}
  \rocq{infotheo.dumas2017dual.dsdp.hopping.dsdp_alice_trace_link.dsdp_run_tracesE,
       infotheo.dumas2017dual.dsdp.hopping.dsdp_alice_trace_link.dsdp_run_traces_encE,
       infotheo.dumas2017dual.dsdp.hopping.dsdp_alice_trace_link.alice_trace_of_hop_tupleE}
  \rocqok
  \uses{def:alice_trace,lem:dsdp_is_correct}
  The interpreter run on the three DSDP party programs terminates with the
  message log the protocol prescribes, and Alice's trace is the image of her
  hop-$0$ tuple under a map $\tau$.
  \[ \mathrm{AliceTrace} = \tau \circ \hoptuple{0} \]
\end{lemma}

\begin{theorem}[The two bounds at the trace]
  \label{thm:alice_trace_bounds}
  \rocq{infotheo.dumas2017dual.dsdp.hopping.dsdp_alice_main.alice_trace_guess_V2_le,
       infotheo.dumas2017dual.dsdp.hopping.dsdp_alice_trace_link.alice_trace_simulator,
       infotheo.dumas2017dual.dsdp.hopping.dsdp_alice_trace_link.alice_trace_ideal,
       infotheo.dumas2017dual.dsdp.hopping.dsdp_alice_main.alice_trace_sim_advantage_le}
  \rocqok
  \uses{lem:trace_of_hop_tuple,def:alice_trace_chain,def:alice_trace_sim_chain,thm:alice_guess_real,thm:alice_simulation_secure}
  With $u_3$ a unit of the plaintext ring, every predictor reading Alice's
  executed trace obeys the guessing bound at $g \circ \tau$. For all weights,
  every distinguisher separates the real trace law from the trace-level ideal
  law by at most the two hop advantages.
  \[ \Pr[\,g(\mathrm{AliceTrace}) = V_2\,]
       \le \tfrac{1}{m} + \epshopsum{D_{g \circ \tau}} \]
  The simulation bound is the gap result of
  Definition~\ref{def:alice_trace_sim_chain}, and the guessing bound is that
  simulation bound at the predictor's distinguisher plus the fiber term of the
  ideal trace, as at the tuple; the program of
  Definition~\ref{def:alice_trace_chain} returns the same bound.
  Proved in \texttt{dsdp\_alice\_main.v}. Both bounds at the raw
  trace, Theorems~\ref{thm:raw_trace_guess}
  and~\ref{thm:raw_trace_simulation}, are read through the decoding of
  Lemma~\ref{lem:raw_trace_decode}.
\end{theorem}

\begin{lemma}[Decoding Alice's trace at her own key]
  \label{lem:raw_trace_decode}
  \rocq{infotheo.dumas2017dual.dsdp.hopping.dsdp_alice_trace_link.alice_raw_trace,
       infotheo.dumas2017dual.dsdp.hopping.dsdp_alice_trace_link.di_data_of_trace_data,
       infotheo.dumas2017dual.dsdp.hopping.dsdp_alice_trace_link.alice_raw_trace_decodeE}
  \rocqok
  \uses{def:alice_trace,lem:trace_of_hop_tuple}
  Alice's raw trace is the message log the interpreter records for her at a
  sample, before the encoding into the finite carrier a distribution is formed
  on. Decoding one encoded entry at a fixed key pair returns a plaintext as
  itself, a ciphertext as itself, and the private-key mark as that key; at
  Alice's own key pair the round trip is exact, her trace carrying no
  public-key mark to constrain.
  \[ \mathrm{map}\ \delta_{\mathrm{dk}_A}\ (\mathrm{AliceTrace}(s))
       = \mathrm{AliceRawTrace}(s) \]
  The encoding therefore takes nothing from Alice and withholds nothing from her: a
  predictor reading either format succeeds on exactly the same samples, so a
  bound at the encoded trace is a bound at the trace the protocol produces.
\end{lemma}

\begin{theorem}[Guessing from the raw interpreter trace]
  \label{thm:raw_trace_guess}
  \rocq{infotheo.dumas2017dual.dsdp.hopping.dsdp_alice_main.alice_raw_trace_guess_V2_le}
  \rocqok
  \uses{thm:alice_trace_bounds,lem:raw_trace_decode}
  With $u_3$ a unit of the plaintext ring, every predictor reading the raw
  trace the interpreter records for Alice when it runs the protocol returns
  Bob's input with probability at most the uniform residue plus the two hop
  advantages of the encoded-trace predictor it induces.
  \[ \Pr[\,g(\mathrm{AliceRawTrace}) = V_2\,]
       \le \tfrac{1}{m} + \epshopsum{D_{g \circ \delta \circ \tau}} \]
  The residue is information-theoretic and each advantage is conditional on the
  IND-CPA assumption at one key. This is the guessing bound at the observation
  a corrupted Alice actually holds, the encoded trace being an intermediate of
  the interpreter rather than something the protocol emits.
  Proved in \texttt{dsdp\_alice\_main.v} as
  \texttt{alice\_raw\_trace\_guess\_V2\_le}.
\end{theorem}

\begin{theorem}[Simulation at the raw trace, averaged over the coin]
  \label{thm:raw_trace_simulation}
  \rocq{infotheo.dumas2017dual.dsdp.hopping.dsdp_alice_main.alice_raw_trace_sim_advantage_avg_le}
  \rocqok
  \uses{thm:alice_trace_bounds,lem:raw_trace_decode}
  With the re-encryption coin drawn uniformly, a Boolean test tells the real
  raw trace from the simulated one, on average over that coin, no more often
  than the average of the two per-coin hop advantages of its decoded lift.
  \[ \big| \Pr[\,D(\mathrm{Real}^{\mathrm{avg}})\,]
       - \Pr[\,D(\mathrm{Ideal}^{\mathrm{avg}})\,] \big|
     \le \textstyle\sum_w \mathcal{U}(w)\,\epshopsum{D_w} \]
  Drawing the coin makes the statement one about the protocol rather than about
  one of its executions, and both terms of the bound are conditional on the
  IND-CPA assumption, one at Bob's key and one at Charlie's.
  Proved in \texttt{dsdp\_alice\_main.v} as
  \texttt{alice\_raw\_trace\_sim\_advantage\_avg\_le}.
\end{theorem}

\begin{theorem}[Trace and tuple leave the same uncertainty]
  \label{thm:centropy_trace_tuple}
  \rocq{infotheo.dumas2017dual.dsdp.hopping.dsdp_alice_trace_link.alice_trace_tupleT,
       infotheo.dumas2017dual.dsdp.hopping.dsdp_alice_trace_link.AliceTraceTuple,
       infotheo.dumas2017dual.dsdp.hopping.dsdp_alice_trace_link.alice_trace_tupleE,
       infotheo.dumas2017dual.dsdp.hopping.dsdp_alice_trace_link.centropy_V2_trace_tupleE}
  \rocqok
  \uses{lem:trace_of_hop_tuple}
  Conditioning on Alice's executed trace leaves the same uncertainty about
  Bob's input as conditioning on her hop-$0$ tuple. The equality holds for all
  weights.
  \[ \HH(V_2 \mid \mathrm{AliceTrace}) = \HH(V_2 \mid \hoptuple{0}) \]
  Proved in \texttt{dsdp\_alice\_trace\_link.v}; with
  Theorem~\ref{thm:centropy_trace_eq0} it puts the zero at the tuple as well.
\end{theorem}

\chapter{The bound conditional on an adversary class}
\label{ch:adversary_class}

The bounds above quantify over every predictor and every distinguisher, and
each hop loses the advantage of one reduction. This chapter adds the two
notions a computational reading of such a bound needs, an asymptotic one and an
adversary class, records what forces the class to leave something out, and
derives the class and its advantage, at Paillier and at Benaloh, from the
$e$-th residuosity assumption.

\begin{definition}[Negligibility of a sequence]
  \label{def:negligible_fun}
  \rocq{infotheo.computational_security.negligible.negligible_fun}
  \rocqok
  A sequence of reals indexed by the security parameter is negligible when it
  eventually falls below every inverse polynomial.
  \[ \mathrm{negl}(f) \iff \forall c,\ \exists N,\ \forall n > N,\
       f(n) < n^{-c} \]
  Every bound in this blueprint is stated at one fixed instance, where an
  asymptotic notion has nothing to measure. This is the shape a sequence of
  instances must have for a bound of that form to vanish faster than every
  inverse polynomial.
\end{definition}

\begin{lemma}[A finite sum of negligible sequences is negligible]
  \label{lem:negligible_fun_sum}
  \rocq{infotheo.computational_security.negligible.negligible_fun_sum}
  \rocqok
  \uses{def:negligible_fun}
  Summing finitely many negligible sequences, indexed by a list, leaves a
  negligible sequence.
  \[ \big(\forall i,\ \mathrm{negl}(F_i)\big) \implies
     \mathrm{negl}\Big(k \mapsto \sum_{i \in s} F_i(k)\Big) \]
  The negligible sequences are an additive submonoid of the sequences of
  reals, and a sum over a list lands in it as soon as each entry does. A
  chain totals one loss term per label of its loss, so a bound read off such
  a total is negligible as soon as each label's loss sequence is.
\end{lemma}

\begin{definition}[An adversary class and the advantage it assumes]
  \label{def:indcpa_assumption}
  \rocq{infotheo.computational_security.indcpa_game.indcpa_epsilon_assumption}
  \rocqok
  \uses{def:indcpa_epsilon}
  An assumption $\mathcal{A}$ is a Boolean classifier on single-query
  real-or-zero adversaries, one advantage $\varepsilon_{\mathrm{cpa}}$, and the
  assertion that every classified adversary stays below it at every key in the
  image of the public-key map.
  \[ \mathcal{A} = \Big(\mathrm{adm},\ \varepsilon_{\mathrm{cpa}},\
       \forall dk,\ \forall D,\ \mathrm{adm}(D) \Rightarrow
       \epsror{\pk(dk)}{D} \le \varepsilon_{\mathrm{cpa}}\Big) \]
  The classifier is extensional. Functional extensionality is in scope and an
  adversary is a record of functions, so two adversaries with the same state
  law, challenge plaintext and decision are one term and receive one Boolean.
  Running time is a property of a syntax the record does not carry.
  The advantage is a parameter at this generic level and derived at the two
  concrete schemes: Definitions~\ref{def:paillier_indcpa_assumption}
  and~\ref{def:benaloh_indcpa_assumption} build a value of this record from a
  residuosity assumption, with $\varepsilon_{\mathrm{cpa}}$ twice the
  residuosity advantage and the bound field proved.
\end{definition}

\begin{definition}[An assumption with a computed class and a proved promise]
  \label{def:cipher_constant_assumption}
  \rocq{infotheo.computational_security.indcpa_game.adv_decide_cipher_constant,
       infotheo.computational_security.indcpa_game.indcpa_epsilon_cipher_constant_eq0,
       infotheo.computational_security.indcpa_game.cipher_constant_assumption}
  \rocqok
  \uses{def:indcpa_assumption}
  The classifier that admits an adversary exactly when its verdict ignores the
  challenge ciphertext is decidable, and every adversary it admits has
  advantage zero.
  \[ \mathrm{adm}(D) \iff \forall c, ch_1, ch_2:\
       D_{\mathrm{dec}}(c, ch_1) = D_{\mathrm{dec}}(c, ch_2),
     \qquad \mathrm{adm}(D) \implies \epsror{\pk}{D} = 0 \]
  Instantiating the record with that classifier and
  $\varepsilon_{\mathrm{cpa}} = 0$
  discharges its bound field by the equality rather than by hypothesis, so it
  is an inhabitant whose promise is proved. Its class is small and the bounds
  conditional on it are correspondingly weak; what it settles is that the record type
  is inhabited with content, and what a classifier that computes looks like.
\end{definition}

\begin{theorem}[The class-conditional trace bound]
  \label{thm:alice_guess_admissible}
  \rocq{infotheo.dumas2017dual.dsdp.hopping.dsdp_alice_main.alice_trace_guess_V2_pr,
       infotheo.dumas2017dual.dsdp.hopping.dsdp_alice_main.alice_trace_guess_V2_admissible_le,
       infotheo.dumas2017dual.dsdp.hopping.dsdp_alice_main.alice_trace_guess_V2_admissible_pq_le}
  \rocqok
  \uses{thm:alice_trace_bounds,def:indcpa_assumption,def:alice_trace_chain_admissible}
  When both reduction adversaries of a trace predictor $g$ are admissible for
  $\mathcal{A}$, and $u_3$ is a unit of the plaintext ring, $g$ matches Bob's
  input with probability at most $1/m$ plus twice the assumed advantage; the
  composite form writes $m$ as $pq$.
  \[ \Pr[\,g(\mathrm{AliceTrace}) = V_2\,]
       \le \tfrac{1}{m} + 2\,\varepsilon_{\mathrm{cpa}} \]
  Three levels of conditionality meet here. The term $1/m$ is
  information-theoretic and
  unconditional, the residue of the leaked output along the solution fiber. The
  term $2\varepsilon_{\mathrm{cpa}}$ is conditional on $\mathcal{A}$, and
  measures the two ciphertext replacements at the advantage $\mathcal{A}$
  assumes rather than at their own values. The two admissibility premises are conditional on
  the class, and are assumed rather than proved.
  The bound is the soundness field of the program of
  Definition~\ref{def:alice_trace_chain_admissible}, which spends the two
  premises at its two hops, so no triangle inequality is taken here.
  Proved in \texttt{dsdp\_alice\_main.v} as
  \texttt{alice\_trace\_guess\_V2\_admissible\_le} and
  \texttt{alice\_trace\_guess\_V2\_admissible\_pq\_le}, and read at each
  scheme in \texttt{dsdp\_instance\_sequence.v}. The composite form takes its
  two primes and the plaintext count from the scheme's own plaintext space
  instead of assuming them.
\end{theorem}

\section{The residuosity assumption and the IND-CPA bound it yields}

\begin{definition}[The $e$-th residuosity problem and an assumption over it]
  \label{def:residuosity_assumption}
  \rocq{infotheo.homomorphic_encryption.residuosity_game.unit_fdist,
       infotheo.homomorphic_encryption.residuosity_game.residue_fdist,
       infotheo.homomorphic_encryption.residuosity_game.residuosity_distinguisher,
       infotheo.homomorphic_encryption.residuosity_game.residuosity_accept,
       infotheo.homomorphic_encryption.residuosity_game.residuosity_advantage,
       infotheo.homomorphic_encryption.residuosity_game.residuosity_assumption,
       infotheo.homomorphic_encryption.residuosity_game.residuosity_admissible_epsilon_leC,
       infotheo.homomorphic_encryption.residuosity_game.decide_constant_assumption}
  \rocqok
  Fix a finite commutative unit ring $T$ and an exponent $e$. The $e$-th
  residuosity problem at $T$ is to tell a uniform unit of $T$ from the $e$-th
  power of a uniform unit; $\mathsf{unit}$ and $\mathsf{res}_e$ name those two
  laws. A distinguisher $D$ is a finite state type, a law on it, and a Boolean
  verdict read off the state and the challenge ring element, and its advantage
  is the absolute gap between its two acceptance probabilities. An assumption
  $\mathcal{B}$ is a Boolean class of distinguishers, one advantage
  $\varepsilon_{\mathrm{res}}$, and the assertion that every classified
  distinguisher stays below it.
  \[ \mathcal{B} = \Big(\mathrm{adm},\ \varepsilon_{\mathrm{res}},\
       \forall D,\ \mathrm{adm}(D) \Rightarrow
       \big\lvert \Pr[D \text{ accepts } \mathsf{res}_e]
         - \Pr[D \text{ accepts } \mathsf{unit}] \big\rvert
       \le \varepsilon_{\mathrm{res}}\Big) \]
  The record has two readings. At $T = \mathbb{Z}_{n^2}$ and $e = n$ it is
  decisional composite residuosity, Paillier 1999, Conjecture 1; at
  $T = \mathbb{Z}_n$ and $e = r$ it is the higher residuosity problem of
  Benaloh 1994. The two schemes of this chapter therefore assume one thing at
  two rings, which is why it is stated once and instantiated twice. Its
  classifier is extensional for the reason given at
  Definition~\ref{def:indcpa_assumption}, so running time stays a property of a
  syntax the record does not carry. \texttt{decide\_constant\_assumption}
  inhabits the record at $\varepsilon_{\mathrm{res}} = 0$, the class of
  distinguishers whose verdict ignores the challenge, whose bound is proved
  rather than assumed. What that settles is that a statement restricted to a
  residuosity class is not empty for want of a record to read it at.
  \texttt{residuosity\_admissible\_epsilon\_leC} states the same bound with the
  two challenge laws exchanged. A hybrid that returns from $\mathsf{unit}$ to
  $\mathsf{res}_e$ spans the gap in that direction, and the absolute value is
  symmetric, so such a step invokes the class once and not twice.
\end{definition}

\begin{lemma}[Multiplication by a unit fixes the uniform unit law]
  \label{lem:unit_fdist_translateE}
  \rocq{infotheo.homomorphic_encryption.residuosity_game.unit_fdist_translateE,
       infotheo.homomorphic_encryption.residuosity_game.unit_fdistmap_translateE,
       infotheo.homomorphic_encryption.residuosity_game.unit_fdistE,
       infotheo.homomorphic_encryption.residuosity_game.residue_fdistE}
  \rocqok
  \uses{def:residuosity_assumption}
  Pushing $\mathsf{unit}$ forward along multiplication by a fixed unit $a$
  returns $\mathsf{unit}$, and the same holds after any Boolean test is applied
  to the result. This is Katz and Lindell 2015, Lemma 11.15.
  \[ (x \mapsto a x)_{*}\,\mathsf{unit} = \mathsf{unit} \]
  Multiplication by $a$ is a bijection of the unit group and a bijection fixes
  a uniform law, so the argument is a permutation argument rather than a
  pointwise computation in the ring. Once a challenge is a uniform unit,
  multiplying it by $g^m$ erases $m$, and that is where a residue-class
  ciphertext hides its plaintext; it is the one step of the reductions below
  that spends no call to the assumption. Read as uniform laws on subsets of
  $T$, $\mathsf{unit}$ is the uniform law on the units of $T$ and
  $\mathsf{res}_e$ the uniform law on the $e$-th residues, which is the
  textbook wording of the problem.
\end{lemma}

\begin{definition}[Decisional composite residuosity at modulus $pq$]
  \label{def:dcr_assumption}
  \rocq{infotheo.computational_security.paillier_indcpa_scheme.dcr_assumption,
       infotheo.computational_security.paillier_indcpa_scheme.dcr_epsilon}
  \rocqok
  \uses{def:residuosity_assumption}
  The assumption above read at the ring $\mathbb{Z}_{(pq)^2}$ and the exponent
  $pq$: a class of distinguishers of a uniform unit of $\mathbb{Z}_{(pq)^2}$
  from a $pq$-th power of one, and one advantage $\varepsilon_{\mathrm{dcr}}$
  they all stay below.
  \[ \mathrm{adm}(D) \Rightarrow \big\lvert
       \Pr[D \text{ accepts } \mathsf{res}_{pq}]
       - \Pr[D \text{ accepts } \mathsf{unit}] \big\rvert
       \le \varepsilon_{\mathrm{dcr}},
     \qquad T = \mathbb{Z}_{(pq)^2} \]
  A $pq$-th power of a uniform unit is the law of a Paillier encryption of
  zero, so a record of this type is Paillier 1999, Conjecture 1 at the modulus
  the scheme runs at. It is the single computational premise every Paillier
  bound of this blueprint is read at.
\end{definition}

\begin{theorem}[Paillier's real-or-zero advantage under composite residuosity]
  \label{thm:paillier_dcr_epsilon_le}
  \rocq{infotheo.computational_security.paillier_indcpa_scheme.paillier_dcr_epsilon_le,
       infotheo.computational_security.paillier_indcpa_scheme.paillier_dcr_admissible_epsilon_le}
  \rocqok
  \uses{def:dcr_assumption,def:residuosity_assumption,lem:unit_fdist_translateE,def:indcpa_epsilon}
  For every private key and every adversary whose two composite-residuosity
  reductions $\mathcal{B}$ admits, the real-or-zero advantage at that key's
  public key is at most twice the advantage $\mathcal{B}$ assumes.
  \[ \epsror{\pk(dk)}{A} \le 2\,\varepsilon_{\mathrm{dcr}} \]
  The two reductions run $A$ against the residuosity challenge. One multiplies
  the challenge by $g^{p_A}$, the key's generator raised to $A$'s own challenge
  plaintext; the other passes the challenge through unchanged. Under
  $\mathsf{res}_{pq}$ they are the real and the zero arm of the IND-CPA
  experiment, by unfolding. The hybrid has three steps and the factor two is
  their loss: one call to $\mathcal{B}$ carries the real arm from
  $\mathsf{res}_{pq}$ to $\mathsf{unit}$, the middle step is free by
  Lemma~\ref{lem:unit_fdist_translateE}, and a second call carries the zero arm
  back. Both summands are conditional on $\mathcal{B}$, so no part of this
  bound is information-theoretic. The single number-theoretic input is
  $g^{pq} = 1$, the order condition the private key record already carries;
  primality of $p$ and $q$ is unused, so the statement is Katz and Lindell
  2015, Theorem 13.13 generalized to any generator whose order divides the
  modulus.
\end{theorem}

\begin{definition}[The IND-CPA assumption of Paillier, derived]
  \label{def:paillier_indcpa_assumption}
  \rocq{infotheo.computational_security.paillier_indcpa_scheme.paillier_indcpa_assumption,
       infotheo.computational_security.paillier_indcpa_scheme.paillier_indcpa_epsilon_le}
  \rocqok
  \uses{thm:paillier_dcr_epsilon_le,def:dcr_assumption,def:indcpa_assumption}
  At Paillier the adversary-class assumption of
  Definition~\ref{def:indcpa_assumption} is a value built from a
  composite-residuosity record rather than a parameter: its class admits the
  adversaries whose two residuosity reductions $\mathcal{B}$ admits, its
  advantage is $2\varepsilon_{\mathrm{dcr}}$, and its promise field is the
  theorem above in place of a hypothesis.
  \[ \mathcal{A} = \big(\text{both reductions admitted by } \mathcal{B},\
       2\,\varepsilon_{\mathrm{dcr}},\
       \text{Theorem~\ref{thm:paillier_dcr_epsilon_le}}\big) \]
  \texttt{paillier\_indcpa\_epsilon\_le} writes that promise out at the
  Paillier ciphertext itself, as the gap between accepting an encryption of the
  chosen plaintext and accepting an encryption of zero, at every private key
  record and with the adversary holding the public key alone. The class is
  inhabited: \texttt{paillier\_dcr\_admissible\_cipher\_constant} admits every
  adversary whose verdict ignores the ciphertext, at the residuosity record of
  advantage zero. Every Paillier bound below reads its
  $\varepsilon_{\mathrm{cpa}}$ off this value, so each is a multiple of the
  composite-residuosity epsilon, at two residuosity calls per key.
\end{definition}

\begin{definition}[Higher residuosity at modulus $n$ and exponent $r$]
  \label{def:benaloh_residuosity_assumption}
  \rocq{infotheo.computational_security.benaloh_indcpa_scheme.benaloh_residuosity_assumption,
       infotheo.computational_security.benaloh_indcpa_scheme.benaloh_residuosity_epsilon}
  \rocqok
  \uses{def:residuosity_assumption}
  The same assumption read at the ring $\mathbb{Z}_n$ and the exponent $r$: a
  class of distinguishers of a uniform unit of $\mathbb{Z}_n$ from an $r$-th
  power of one, and one advantage $\varepsilon_{\mathrm{res}}$ they all stay
  below.
  \[ \mathrm{adm}(D) \Rightarrow \big\lvert
       \Pr[D \text{ accepts } \mathsf{res}_r]
       - \Pr[D \text{ accepts } \mathsf{unit}] \big\rvert
       \le \varepsilon_{\mathrm{res}},
     \qquad T = \mathbb{Z}_n \]
  This is Benaloh 1994's higher residuosity assumption, and an $r$-th power of
  a uniform unit is the law of a Benaloh encryption of zero. The ring is the
  ciphertext modulus $n$ and the exponent is the block size $r$, which is the
  plaintext count; the two numbers are unrelated.
\end{definition}

\begin{theorem}[Benaloh's real-or-zero advantage under higher residuosity]
  \label{thm:benaloh_residuosity_epsilon_le}
  \rocq{infotheo.computational_security.benaloh_indcpa_scheme.benaloh_residuosity_epsilon_le,
       infotheo.computational_security.benaloh_indcpa_scheme.benaloh_residuosity_admissible_epsilon_le}
  \rocqok
  \uses{def:benaloh_residuosity_assumption,def:residuosity_assumption,lem:unit_fdist_translateE,def:indcpa_epsilon}
  For every private key and every adversary whose two residuosity reductions
  $\mathcal{B}$ admits, the real-or-zero advantage at that key's public key is
  at most twice the advantage $\mathcal{B}$ assumes.
  \[ \epsror{\pk(dk)}{A} \le 2\,\varepsilon_{\mathrm{res}} \]
  The hybrid is the one of Theorem~\ref{thm:paillier_dcr_epsilon_le} with two
  differences. The multiplier is $y^{p_A}$ at the key's generator $y$, a unit
  of $\mathbb{Z}_n$ by its type, so the order condition $y^r = 1$ the key
  record carries is never read and the bound holds at any generator. And the
  class quantifies over the unit group of $\mathbb{Z}_n$ rather than over the
  ring, that group being where a Benaloh generator lives, which is a Boolean
  quantifier because the group is a finite type. Beyond $1 < n$ and $1 < r$ no
  number theory enters, primality of the factors of $n$ and the structure of
  the unit group included.
\end{theorem}

\begin{definition}[The IND-CPA assumption of Benaloh, derived]
  \label{def:benaloh_indcpa_assumption}
  \rocq{infotheo.computational_security.benaloh_indcpa_scheme.benaloh_indcpa_assumption,
       infotheo.computational_security.benaloh_indcpa_scheme.benaloh_indcpa_epsilon_le}
  \rocqok
  \uses{thm:benaloh_residuosity_epsilon_le,def:benaloh_residuosity_assumption,def:indcpa_assumption}
  At Benaloh the adversary-class assumption is likewise a value built from a
  residuosity record: its class admits the adversaries whose two residuosity
  reductions $\mathcal{B}$ admits, its advantage is
  $2\varepsilon_{\mathrm{res}}$, and its promise field is the theorem above.
  \[ \mathcal{A} = \big(\text{both reductions admitted by } \mathcal{B},\
       2\,\varepsilon_{\mathrm{res}},\
       \text{Theorem~\ref{thm:benaloh_residuosity_epsilon_le}}\big) \]
  \texttt{benaloh\_indcpa\_epsilon\_le} writes that promise out at the Benaloh
  ciphertext, and
  \texttt{benaloh\_residuosity\_admissible\_cipher\_constant} inhabits the
  class at the residuosity record of advantage zero.
\end{definition}

\begin{corollary}[The bound at a Paillier instance]
  \label{cor:alice_guess_paillier}
  \rocq{infotheo.dumas2017dual.dsdp.hopping.dsdp_alice_instances.paillier_trace_guess_V2_admissible_pq_le,
       infotheo.dumas2017dual.dsdp.hopping.dsdp_alice_instances.paillier_epsilon_at_dcrE}
  \rocqok
  \uses{thm:alice_guess_admissible,def:paillier_setting,def:dcr_assumption,def:paillier_indcpa_assumption}
  At one Paillier instance of modulus $p q$, a trace predictor whose two
  reduction adversaries the class induced by $\mathcal{B}$ admits matches
  Bob's input with probability at most $1/(p q)$ plus four times the assumed
  residuosity advantage.
  \[ \Pr[\,g(\mathrm{AliceTrace}) = V_2\,]
       \le \tfrac{1}{p q} + 4\,\varepsilon_{\mathrm{dcr}} \]
  Both sides are discharged at this scheme. The information-theoretic side
  reads the plaintext count off the packaging, whose plaintext space has
  exactly $p q$ elements, and the coin plumbing the abstract bound quantifies
  over is supplied by the scheme's own randomness. The computational side is
  decisional composite residuosity and nothing further, at the loss
  Definition~\ref{def:paillier_indcpa_assumption} fixes: the factor four is
  two keys, Bob's and Charlie's, times two residuosity calls per key, the
  identity \texttt{paillier\_epsilon\_at\_dcrE}. What is assumed here is
  $\mathcal{B}$, and the class premises restrict which predictors the bound
  covers. The bound holds at each fixed modulus on its own, with no growth
  condition on $p$ and $q$. The sequence statements below read it along a
  sequence of such instances.
\end{corollary}

\begin{corollary}[The bound at a Benaloh instance]
  \label{cor:alice_guess_benaloh}
  \rocq{infotheo.dumas2017dual.dsdp.hopping.dsdp_alice_instances.benaloh_trace_guess_V2_admissible_pq_le,
       infotheo.dumas2017dual.dsdp.hopping.dsdp_alice_instances.benaloh_epsilon_at_residuosityE}
  \rocqok
  \uses{thm:alice_guess_admissible,def:benaloh_setting,def:benaloh_residuosity_assumption,def:benaloh_indcpa_assumption}
  At one Benaloh instance whose plaintext space has $r = p q$ elements, a
  trace predictor whose two reduction adversaries the class induced by
  $\mathcal{B}$ admits matches Bob's input with probability at most
  $1/(p q)$ plus four times the assumed residuosity advantage.
  \[ \Pr[\,g(\mathrm{AliceTrace}) = V_2\,]
       \le \tfrac{1}{p q} + 4\,\varepsilon_{\mathrm{res}} \]
  The block size is where the unconditional term is read, so the factored
  form $r = p q$ enters as a premise of this reading and not of the scheme,
  which asks only $1 < r$. The computational side is $r$-th residuosity at
  the loss Definition~\ref{def:benaloh_indcpa_assumption} fixes, the factor
  four again two keys times two residuosity calls per key, the identity
  \texttt{benaloh\_epsilon\_at\_residuosityE}.
\end{corollary}

\section{epsHop, the language a multi-hop bound is written in}

Each of the three reductions above is a hybrid: a sequence of games joined by
steps, some of which spend an assumption and some of which are identities. This
section states the language they are written in, epsHop. A chain is a morphism
between acceptance probabilities graded by an accumulated loss, and that loss is
a list of labels rather than a single real, so a finished derivation carries the
assumptions its bound rests on alongside their numeric total. epsHop
is the dual of the piSMC language of \texttt{smc/pismc.v}: piSMC writes what the
parties of a protocol do, epsHop writes what a security argument about them
loses. The three programs of the section are the Paillier hybrid, the Benaloh
hybrid, and Alice's two-hop ladder, each written in it, and the section closes
with the reading of such a program along a sequence of security parameters.

A program is one delimited block, $\varepsilon\{\;\cdots\;\}$, whose statements
are separated by a semicolon and read as follows.
\begin{itemize}
  \item \texttt{start g} opens the chain at the game $g$, with no loss yet.
  \item \texttt{hop l e to g' by H} invokes the assumption named $l$, loses
    $e$, and reaches the game $g'$, guaranteed by
    $H : \lvert \mathrm{current} - g' \rvert \le e$. The label $l$ fixes the
    source game, the target and the loss through its claim, and $e$, $g'$ and
    $H$ are each checked against it.
  \item \texttt{same to g' by H} rewrites the game to $g'$ at no loss,
    guaranteed by $H : \mathrm{current} = g'$.
  \item \texttt{plus l c by H} is \texttt{hop l c to 0 by H}, under a label
    whose target is the zero game: the last game lies within $c$ of zero,
    guaranteed by $H : \lvert \mathrm{current} - 0 \rvert \le c$, so it is at
    most $c$ and $c$ enters the loss. It is spelled \texttt{plus} so that a
    reader sees a term added rather than a hop taken.
  \item \texttt{s ;; bound c by H} returns the bound: the total of the loss is
    $c$, guaranteed by $H : \mathrm{total} = c$, so the advantage of the
    chain, the distance between the game it opened at and the game it stopped
    at, is at most $c$.
\end{itemize}
The label slot names the assumption invoked and where, the loss slot is that
assumption's advantage, and the proof slot says whether the term is assumed or
exact. All three are compared with what the label claims. Each proof slot is
mandatory: a statement is ill-formed without it, and a proof about the wrong
pair of games neither convinces nor names a step. So a chain that exists has no
pending obligation, and a bound is read off it by projection.

\begin{definition}[A label and the claim it stands for]
  \label{def:claim}
  \rocq{infotheo.computational_security.epshop.claim,
       infotheo.computational_security.epshop.hop_obligation}
  \rocqok
  A label is an element of a type the client chooses, and a claim function
  $C$ sends it to what it asserts. A claim $\mathrm{Claim}(s, t, e)$ names a
  source game, a target game and a loss term, and asserts that the two games
  lie within that loss of each other.
  \[ \mathrm{obl}(\mathrm{Claim}(s,t,e))
       = \big(\lvert s - t \rvert \le e\big) \]
  The three components are read off by the projections $\mathrm{src}$,
  $\mathrm{tgt}$ and $\mathrm{loss}$, so the loss term of a label $l$ is
  $\mathrm{loss}(C(l))$. The claim is the whole content of a
  label: it is what a step written under that label is checked against, so a
  step cannot record a call to an assumption it did not make. A label whose
  target is the zero game is what a bound on one game is written under, that
  game's distance from zero being the game itself.
\end{definition}

\begin{definition}[A loss and the labels it is a list of]
  \label{def:loss}
  \rocq{infotheo.computational_security.epshop.loss,
       infotheo.computational_security.epshop.loss_eval,
       infotheo.computational_security.epshop.loss_eval_nil,
       infotheo.computational_security.epshop.loss_eval1,
       infotheo.computational_security.epshop.loss_eval_cat,
       infotheo.computational_security.epshop.loss_total,
       infotheo.computational_security.epshop.loss_evalE}
  \rocqok
  \uses{def:claim}
  A loss is a list of labels, the free monoid on them, and its total
  $\mathrm{ev}$ sums the loss term each label's claim names, a monoid morphism to
  the additive reals.
  \[ \mathrm{ev}(s) = \sum_{l \in s} \mathrm{loss}(C(l)), \qquad
     \mathrm{ev}[\,] = 0, \qquad \mathrm{ev}[l] = \mathrm{loss}(C(l)) \]
  \[ \mathrm{ev}(s_1 \mathbin{+\!\!+} s_2)
       = \mathrm{ev}(s_1) + \mathrm{ev}(s_2) \]
  The same total is also written as a left fold seeded at the first label's
  loss term, and the two agree. On a list of literal labels the fold reduces to
  a left-associated sum of the loss terms those labels claim, with no summation and
  no list left in it.
  \[ \mathrm{tot}(s) = \mathrm{ev}(s), \qquad
     \mathrm{tot}[l_1, l_2, l_3]
       = \big(\mathrm{loss}(C(l_1)) + \mathrm{loss}(C(l_2))\big)
         + \mathrm{loss}(C(l_3)) \]
  A label names the assumption and the reduction that one step invokes, and the
  total forgets it. The two readings answer different questions about the same
  bound: the list says of each summand whether it is information-theoretic or
  conditional on a computational assumption and at which key, the total is the
  number a theorem states. That $\mathrm{ev}$ is a morphism is what makes the
  total of a composite the sum of the totals of its parts.
\end{definition}

\begin{definition}[A chain and the four statements that build one]
  \label{def:chain}
  \rocq{infotheo.computational_security.epshop.chain,
       infotheo.computational_security.epshop.chain_start,
       infotheo.computational_security.epshop.chain_hop,
       infotheo.computational_security.epshop.chain_same,
       infotheo.computational_security.epshop.chain_then,
       infotheo.computational_security.epshop.le_of_eq,
       infotheo.computational_security.epshop.plus_le,
       infotheo.computational_security.epshop.advantage0}
  \rocqok
  \uses{def:loss}
  A chain is a head endpoint $a$, a last endpoint $b$, a loss $s$, and a
  proof that the total of the loss bounds the distance between the endpoints.
  The endpoints are acceptance probabilities.
  \[ \lvert a - b \rvert \le \mathrm{ev}(s) \]
  Four statements build one, each re-establishing that invariant and each
  taking its own justification as an argument, and are written below as the
  three data $(a, b, s)$ they carry.
  \[ \mathrm{start}(g) = \big(g,\; g,\; [\,]\big), \qquad
     \mathrm{hop}(l, e, g', H) = \big(\mathrm{src}(C(l)),\; g',\; [l]\big)
       \quad\text{for } H : \mathrm{obl}(C(l)) \]
  \[ \mathrm{same}(g', H) = \big(x,\; g',\; [\,]\big)
       \quad\text{for } H : x = g', \qquad
     \mathrm{then}(m, f, H_b) = \big(a_m,\; b_f,\;
       s_m \mathbin{+\!\!+} s_f\big) \]
  An assumption enters a derivation at $\mathrm{hop}$ and nowhere else, which
  is why the loss of a finished chain is the complete list of what its bound
  rests on. A $\mathrm{hop}$ takes two further arguments, the equations
  $e = \mathrm{loss}(C(l))$ and $g' = \mathrm{tgt}(C(l))$, each discharged by
  reflexivity, so the loss and the target written at a step are the loss and
  the target the label claims. $\mathrm{same}$ is the step an
  information-theoretic identity between two games takes, and it leaves the
  loss as it stands; \texttt{le\_of\_eq} turns such an identity into the
  inequality a $\mathrm{hop}$ asks for, which is what a hop whose gap is an
  equality of advantages rather than an assumed bound is justified by, and
  \texttt{plus\_le} does the same for a hop whose target is the zero game,
  \texttt{advantage0} reading the distance from zero back as the game. In
  $\mathrm{hop}$ the source endpoint is the source of the label's claim and in
  $\mathrm{same}$ it is read off the type of the justification, so neither
  writes it. Composition $\mathrm{then}$ is partial, sound
  only under the side condition $H_b : a_f = b_m$, that the second fragment
  starts where the first stopped, and it is where the one triangle inequality
  of the development is discharged, once, so that a chain of any length needs
  none. Without that side condition the concatenated loss would bound nothing
  about the two outer endpoints. The semicolon of a program discharges it by
  conversion, so two consecutive statements type-check exactly when the second
  opens at the game the first reached.
\end{definition}

\begin{definition}[What a program returns]
  \label{def:chain_result}
  \rocq{infotheo.computational_security.epshop.chain_result,
       infotheo.computational_security.epshop.chain_result_of_chain,
       infotheo.computational_security.epshop.chain_bound}
  \rocqok
  \uses{def:chain}
  A result is an advantage $a$, a loss $s$, a bound $c$, a proof that $c$
  bounds $a$, and a proof that $c$ is the total of $s$.
  \[ a \le c, \qquad \mathrm{tot}(s) = c \]
  A chain returns one on its own, on the distance between its endpoints, at
  the total of its loss; $\mathrm{bound}$ republishes a result at an explicit
  $c$ for $H : \mathrm{tot}(s) = c$.
  \[ m = \big(\lvert a_m - b_m \rvert,\; s_m,\; \mathrm{tot}(s_m)\big),
     \qquad
     \mathrm{bound}(b, c, H) = \big(a_b,\; s_b,\; c\big) \]
  A chain whose last game is the zero game therefore returns a bound on the
  probability of the game it opened at, which is the shape a secrecy theorem
  has, and the label of that last step is the one summand of such a bound that
  comes from the endpoint rather than from a ciphertext replacement.
  $\mathrm{bound}$ is where the published number is written
  down: since $\mathrm{tot}$ converts to a plain sum on a list of literal
  labels, $H$ is an algebraic identity between that sum and the number the
  client's theorem states, and no vocabulary of the language appears in it.
  It is written after the double semicolon because it returns a result
  where every other statement returns a chain. The last field is what keeps
  the label list, rather than the number written at $\mathrm{bound}$, the
  source of a published bound: republishing composes $H$ with the identity
  the result already carried, so the bound of a program is the total of its
  labels however often it has been rewritten. A reading of a bound along a
  sequence of security parameters is a reading of that list, and it is sound
  only because of this field.
\end{definition}

\begin{lemma}[The category laws, as equalities of records]
  \label{lem:chain_laws}
  \rocq{infotheo.computational_security.epshop.chain_observable_eq,
       infotheo.computational_security.epshop.chain_left_unit,
       infotheo.computational_security.epshop.chain_right_unit,
       infotheo.computational_security.epshop.chain_assoc}
  \rocqok
  \uses{def:chain}
  Two chains that agree on their three observable fields are equal, and the
  unit and associativity laws hold as equalities of records.
  \[ \mathrm{then}(\mathrm{start}(g), \mathrm{frag}) = \mathrm{frag}, \qquad
     \mathrm{then}(m, \mathrm{start}(b_m)) = m \]
  \[ \mathrm{then}(\mathrm{then}(m_1, m_2), m_3)
       = \mathrm{then}(m_1, \mathrm{then}(m_2, m_3)) \]
  The fourth field of a chain is a proof of a Boolean, whose uniqueness is a
  lemma of the library and not an axiom, so \texttt{chain\_observable\_eq}
  uses nothing classical and the three laws are equalities rather than
  statements up to an equivalence. Their content for a program is that the
  regrouping and unit rearrangements the semicolon performs leave the loss, and
  therefore every bound of the three programs below, as it stands.
\end{lemma}

\begin{definition}[The Paillier hybrid as a chain]
  \label{def:paillier_chain}
  \rocq{infotheo.computational_security.paillier_indcpa_scheme.paillier_label,
       infotheo.computational_security.paillier_indcpa_scheme.paillier_claim,
       infotheo.computational_security.paillier_indcpa_scheme.paillier_chain,
       infotheo.computational_security.paillier_indcpa_scheme.dcr_totalE}
  \rocqok
  \uses{def:chain,def:chain_result,def:claim,def:loss,def:dcr_assumption,thm:paillier_dcr_epsilon_le,lem:unit_fdist_translateE}
  The three-step hybrid of Theorem~\ref{thm:paillier_dcr_epsilon_le} as a chain
  over the four acceptance probabilities it passes through. It starts at the
  multiplying reduction under $\mathsf{res}_{pq}$, which is the real arm of the
  IND-CPA experiment; hops to that reduction under $\mathsf{unit}$; crosses to
  the plain reduction under $\mathsf{unit}$ by
  Lemma~\ref{lem:unit_fdist_translateE}, which is an identity and logs nothing;
  and hops to the plain reduction under $\mathsf{res}_{pq}$, which is the zero
  arm. Its loss is the two labels $\mathsf{dcr\_g}$ and $\mathsf{dcr\_0}$,
  named after the reduction each hop invokes the assumption through, and
  \texttt{paillier\_claim} fixes for each of them the two experiments its call
  moves between and the epsilon it assumes.
  It returns the bound $\mathrm{tot}(s)$ in closed form, the identity
  \texttt{dcr\_totalE} proving the return statement.
  \[ \mathrm{tot}(s) = \varepsilon_{\mathrm{dcr}} + \varepsilon_{\mathrm{dcr}}
       = 2\,\varepsilon_{\mathrm{dcr}} \]
  The two class memberships $\mathrm{adm}(\mathcal{B}_{\mathrm{mul}})$ and
  $\mathrm{adm}(\mathcal{B}_{\mathrm{pl}})$ are parameters of the chain, spent
  at the two hops as they are built. Both terms are conditional on
  $\mathcal{B}$, so this loss has one kind of term. The number in
  Theorem~\ref{thm:paillier_dcr_epsilon_le} is the bound this program returns
  and its two hypotheses are those two memberships, so that theorem is the
  soundness field of this result rather than a triangle inequality of its own.
\end{definition}

\begin{definition}[The Benaloh hybrid as a chain]
  \label{def:benaloh_chain}
  \rocq{infotheo.computational_security.benaloh_indcpa_scheme.benaloh_label,
       infotheo.computational_security.benaloh_indcpa_scheme.benaloh_claim,
       infotheo.computational_security.benaloh_indcpa_scheme.benaloh_chain,
       infotheo.computational_security.benaloh_indcpa_scheme.residuosity_totalE}
  \rocqok
  \uses{def:chain,def:chain_result,def:claim,def:loss,def:benaloh_residuosity_assumption,thm:benaloh_residuosity_epsilon_le,lem:unit_fdist_translateE}
  The same three steps at $T = \mathbb{Z}_n$ and the exponent $r$, with the
  labels $\mathsf{residuosity\_y}$ and $\mathsf{residuosity\_0}$, their claims
  fixed by \texttt{benaloh\_claim}, and the same two-label loss.
  \[ \mathrm{tot}(s) = \varepsilon_{\mathrm{res}} + \varepsilon_{\mathrm{res}}
       = 2\,\varepsilon_{\mathrm{res}} \]
  One place separates it from the Paillier chain. Its middle identity takes no
  order premise on the generator, the multiplier $y^{p_A}$ being a unit of
  $\mathbb{Z}_n$ by its type, so the $\mathrm{same}$ step there is discharged
  without reading the key's order condition. Both terms are conditional on
  $\mathcal{B}$, and Theorem~\ref{thm:benaloh_residuosity_epsilon_le} is again
  the soundness field of this result at its two class memberships.
\end{definition}

\begin{definition}[Alice's trace argument as a chain opening at the protocol]
  \label{def:alice_trace_chain}
  \rocq{infotheo.dumas2017dual.dsdp.hopping.dsdp_alice_main.alice_label,
       infotheo.dumas2017dual.dsdp.hopping.dsdp_alice_main.alice_claim,
       infotheo.dumas2017dual.dsdp.hopping.dsdp_alice_main.all_zero_game_V2_le_invm,
       infotheo.dumas2017dual.dsdp.hopping.dsdp_alice_main.alice_totalE,
       infotheo.dumas2017dual.dsdp.hopping.dsdp_alice_main.accept_trace_tupleE,
       infotheo.dumas2017dual.dsdp.hopping.dsdp_alice_main.hop_tuple_distinguisher,
       infotheo.dumas2017dual.dsdp.hopping.dsdp_alice_main.alice_trace_chain}
  \rocqok
  \uses{def:chain,def:chain_result,def:claim,def:loss,def:alice_trace,lem:trace_of_hop_tuple,lem:alice_view_advantage,lem:guess_all_zero}
  Alice's guessing argument as one chain, opening at the executed protocol.
  $G_\tau$ is the law of the honest input pair beside the trace the interpreter
  returns for Alice when it runs $\mathrm{dsdp\_protocol}$ at a sample, and
  $\gamma_\tau(D_g)$ the probability that a predictor's distinguisher accepts a
  value sampled from it, so the object the argument starts from is the executed
  protocol and not the tuple $\hoptuple{0}$.
  \[ \mathrm{start}\ \gamma_\tau(D_g) \]
  \[ \mathrm{same\ to}\ \gamma_0(D_{g \circ \tau})\ \mathrm{by}\
       \texttt{accept\_trace\_tupleE} \]
  The step to $\gamma_0$ carries no label and adds nothing to the loss: by
  Lemma~\ref{lem:trace_of_hop_tuple} the trace is the image of the tuple under
  $\tau$, so a test on the trace accepts as often as its lift
  $D_{g \circ \tau}$ on the tuple, which is
  \texttt{hop\_tuple\_distinguisher} applied once under a local name. Two hops
  then run the ciphertext replacements over $\gamma_0$, $\gamma_1$ and
  $\gamma_2$ at that lift, each logging the advantage of the reduction its
  label names, at Bob's key and then at Charlie's, and a third step hops from
  the all-zero endpoint to the zero game, so the program is a chain bound on
  $\gamma_\tau(D_g)$. \texttt{alice\_claim} fixes what each of the three labels
  claims: the games a hop label joins and the advantage it loses, and the
  all-zero game against the zero game.
  \[ \mathrm{tot}(s) = \tfrac{1}{m} + \epshopsum{D_{g \circ \tau}} \]
  The chain accumulates the two hop terms before the endpoint term, so the
  return statement reorders, which is what \texttt{alice\_totalE} proves.
  This loss has two kinds of term and the labels separate them. The two hop
  terms are the advantages of the reductions $\hopred{0}{D_{g \circ \tau}}$ and
  $\hopred{1}{D_{g \circ \tau}}$, each of which becomes conditional on a
  computational assumption only where a later statement requires it to be
  small. The term $\mathsf{uniform\_fiber}$ is information-theoretic and
  unconditional, the residue the leaked output leaves along the solution fiber.
  Each hop is justified by Lemma~\ref{lem:alice_view_advantage}, which makes
  its gap exactly the term it logs, read as an inequality through
  \texttt{le\_of\_eq}; the step to the zero game is justified by
  \texttt{all\_zero\_game\_V2\_le\_invm}, which is
  Lemma~\ref{lem:guess_all_zero} carried to the game form and is the one place
  the fiber enters. The trace bound of Theorem~\ref{thm:alice_trace_bounds} is
  the soundness field of this result.
\end{definition}

\begin{definition}[Alice's trace against the simulation as one chain]
  \label{def:alice_trace_sim_chain}
  \rocq{infotheo.dumas2017dual.dsdp.hopping.dsdp_alice_main.accept_trace_ideal_tupleE,
       infotheo.dumas2017dual.dsdp.hopping.dsdp_alice_main.alice_trace_sim_chain}
  \rocqok
  \uses{def:alice_trace_chain,def:alice_trace,lem:trace_of_hop_tuple,lem:simulator_factorization}
  The two hops of Definition~\ref{def:alice_trace_chain} with one step at each
  end that adds nothing to the loss. The chain opens at $\gamma_\tau(D)$, runs the
  two ciphertext replacements at the lifted test $D \circ \tau$, and closes at
  $\Pr_{\Ideal^\tau}[D]$, the probability that the same test accepts the
  trace-level ideal law.
  \[ \mathrm{same\ to}\ \gamma_0(D \circ \tau)\ \mathrm{by}\
       \texttt{accept\_trace\_tupleE},
     \qquad
     \mathrm{same\ to}\ \Pr_{\Ideal^\tau}[D]\ \mathrm{by}\
       \texttt{accept\_trace\_ideal\_tupleE} \]
  The first step is Lemma~\ref{lem:trace_of_hop_tuple} read at acceptance
  probabilities. The second holds because $\Ideal^\tau$ is the image under
  $\tau$ of $\Ideal$, which Lemma~\ref{lem:simulator_factorization} identifies
  with the law of the honest inputs beside $\hoptuple{2}$.
  \[ \mathrm{loss}(s) = [\mathsf{cpa\_bob},\ \mathsf{cpa\_charlie}],
     \qquad \mathrm{tot}(s) = \epshopsum{D \circ \tau} \]
  The chain carries no terminal statement, so the gap result it returns is the
  trace-level simulation bound of Theorem~\ref{thm:alice_trace_bounds}.
  \[ \big\lvert \gamma_\tau(D) - \Pr_{\Ideal^\tau}[D] \big\rvert
       \le \epshopsum{D \circ \tau} \]
\end{definition}

\begin{definition}[Alice's trace argument at the class epsilon]
  \label{def:alice_trace_chain_admissible}
  \rocq{infotheo.dumas2017dual.dsdp.hopping.dsdp_alice_main.alice_claim_admissible,
       infotheo.dumas2017dual.dsdp.hopping.dsdp_alice_main.alice_admissible_totalE,
       infotheo.dumas2017dual.dsdp.hopping.dsdp_alice_main.f_guess_V2_le}
  \rocqok
  \uses{def:alice_trace_chain,def:indcpa_assumption,def:claim,def:loss}
  The program of Definition~\ref{def:alice_trace_chain} over the dictionary
  $C^{\mathcal{A}}$ that charges each hop the advantage $\mathcal{A}$ assumes
  in place of the advantage its own reduction shows. The three games and the
  three labels are those of \texttt{alice\_claim}.
  \[ C^{\mathcal{A}}(\mathsf{cpa\_bob})
       = \big(\gamma_0(D),\ \gamma_1(D),\ \varepsilon_{\mathrm{cpa}}\big),
     \qquad
     C^{\mathcal{A}}(\mathsf{cpa\_charlie})
       = \big(\gamma_1(D),\ \gamma_2(D),\ \varepsilon_{\mathrm{cpa}}\big) \]
  The two class memberships are premises of the statements the program is
  written under, and are spent in the justification of the hop each one
  licenses, which is the shape Definition~\ref{def:paillier_chain} has. The zero-game label loses the same
  unconditional residue as in Definition~\ref{def:alice_trace_chain}, so the
  total adds one unconditional term to two class-conditional ones, and
  \texttt{alice\_admissible\_totalE} is the reordering that states it.
  \[ \mathrm{tot}(s) = \tfrac{1}{m} + 2\,\varepsilon_{\mathrm{cpa}} \]
  The program is written along a sequence of instances, under
  Theorem~\ref{thm:alice_guess_negligible} and again under the bound it
  returns at each parameter, and Theorem~\ref{thm:alice_guess_admissible} is
  that bound read at the sequence which repeats one instance.
\end{definition}

\begin{definition}[Alice's trace against the simulation at the class epsilon]
  \label{def:alice_trace_sim_chain_admissible}
  \rocq{infotheo.dumas2017dual.dsdp.hopping.dsdp_alice_main.alice_trace_sim_advantage,
       infotheo.dumas2017dual.dsdp.hopping.dsdp_alice_main.f_sim_advantage,
       infotheo.dumas2017dual.dsdp.hopping.dsdp_alice_main.f_sim_advantageE}
  \rocqok
  \uses{def:alice_trace_sim_chain,def:alice_trace_chain_admissible}
  The chain of Definition~\ref{def:alice_trace_sim_chain} over the dictionary
  $C^{\mathcal{A}}$ of Definition~\ref{def:alice_trace_chain_admissible}, which
  charges each of the two hops the advantage $\mathcal{A}$ assumes in place of
  the advantage its own reduction shows. The two class memberships are
  premises of the statement the program is written under, and are spent in the
  justification of the hop each one licenses.
  \[ \mathrm{loss}(s) = [\mathsf{cpa\_bob},\ \mathsf{cpa\_charlie}],
     \qquad \mathrm{tot}(s) = 2\,\varepsilon_{\mathrm{cpa}} \]
  The two steps at the ends carry the test between the executed trace and the
  hopping tuple at no loss, and the chain carries no terminal statement, so
  the quantity it opens at is the distance a trace test sees between the
  executed trace and the simulation, and the gap result it returns bounds that
  distance by two class-conditional terms and nothing else.
  \[ \mathrm{adv}^\tau(D)
       = \big\lvert \gamma_\tau(D) - \Pr_{\Ideal^\tau}[D] \big\rvert
       = a_s, \qquad
     \mathrm{adv}^\tau(D) \le 2\,\varepsilon_{\mathrm{cpa}} \]
\end{definition}

Each of the three programs states a bound at one security parameter, where an
asymptotic notion has nothing to measure. The three statements that close this
section read a program along a sequence of parameters instead, which is the
form a computational claim finally takes. A program indexed by the parameter is
a sequence of results, and negligibility is a property of that sequence that no
single member of it can carry.

\begin{definition}[A dictionary all of whose labels lose a negligible sequence]
  \label{def:negligible_claims}
  \rocq{infotheo.computational_security.epshop_sequence.negligibleClaims}
  \rocqok
  \uses{def:claim,def:negligible_fun}
  A sequence of dictionaries $C_k$ indexed by the security parameter, together
  with the negligibility of the loss term of every label of the label type, read
  along that parameter.
  \[ \forall l,\ \mathrm{negl}\big(k \mapsto \mathrm{loss}(C_k(l))\big) \]
  The condition quantifies over the label type and not over the labels one
  program spends, so whatever a chain over $C$ names is covered by it. It is
  attached to the dictionary once, and a theorem about any program written at
  that dictionary then names no label: the dictionary carries the per-label
  statements the theorem would otherwise have to list. Recovering the
  statements from a program is by the name of its dictionary, so a dictionary
  written as an anonymous function can be given this data and never found
  again.
\end{definition}

\begin{lemma}[The total of a label list along the security parameter]
  \label{lem:loss_eval_negligible}
  \rocq{infotheo.computational_security.epshop_sequence.loss_eval_negligible}
  \rocqok
  \uses{def:negligible_claims,def:loss,lem:negligible_fun_sum}
  The total of any list of labels of such a dictionary is a negligible
  sequence.
  \[ \mathrm{negl}\big(k \mapsto \mathrm{ev}_{C_k}(s)\big) \]
  $\mathrm{ev}$ is a monoid morphism out of the free monoid on labels and the
  negligible sequences are a submonoid of the sequences of reals, so a total
  lands in that submonoid exactly when each of its generators does. Nothing
  about membership in $s$ is required, the dictionary covering the whole label
  type.
\end{lemma}

\begin{theorem}[A sequence of programs has a negligible advantage]
  \label{thm:advantage_negligible}
  \rocq{infotheo.computational_security.epshop_sequence.advantage_negligible}
  \rocqok
  \uses{def:negligible_claims,def:chain_result,def:negligible_fun,lem:loss_eval_negligible}
  Let $P$ be a program at every security parameter, over a dictionary all of
  whose labels are negligible, and let its loss be the same at every parameter.
  Then any $f$ identified with the advantage $P$ bounds is negligible.
  \[ \big(\forall k,\ s_{P_k} = s_{P_0}\big) \implies
     \big(\forall k,\ f(k) = a_{P_k}\big) \implies \mathrm{negl}(f) \]
  The program supplies the two facts the argument runs on: its advantage is
  below its published bound, and that bound is the total of its labels. So the
  asymptotic reading of a multi-hop bound is a reading of the label list, and
  the hypothesis that the loss does not vary with the parameter is what makes
  it one list. A bound stated at a fixed instance cannot state this, and the
  step from a theorem of the previous sections to its asymptotic form is
  exactly this statement.
\end{theorem}

\section{What the class leaves out, and the three instance sequences}

\begin{definition}[The decrypting predictor]
  \label{def:trace_decryptor}
  \rocq{infotheo.dumas2017dual.dsdp.hopping.dsdp_alice_trace_link.bob_decrypt_predictor,
       infotheo.dumas2017dual.dsdp.hopping.dsdp_alice_trace_link.alice_trace_decode_V2E}
  \rocqok
  \uses{def:alice_trace}
  The predictor $\delta$ reads Bob's ciphertext off slot $3$ of Alice's trace
  and decrypts it with Bob's private key. Correctness of decryption makes Bob's
  input a function of the trace.
  \[ \delta(t) = \Dec(dk_B,\, t_3),\qquad V_2 = \delta \circ \mathrm{AliceTrace} \]
  It holds a private key, so it lies outside the attack model the hop ladder
  bounds, which grants the public keys alone.
\end{definition}

\begin{theorem}[The class cannot admit every adversary]
  \label{thm:decrypt_lower_bound}
  \rocq{infotheo.dumas2017dual.dsdp.hopping.dsdp_alice_main.decrypt_bob_epsilon_ge}
  \rocqok
  \uses{def:trace_decryptor,thm:alice_guess_real,lem:guess_all_zero}
  With $u_3$ a unit of the plaintext ring, the hop-$0$ advantage the decrypting
  predictor buys is at least $1 - 1/m$.
  \[ 1 - \tfrac{1}{m} \le \epsror{\pk_B}{\hopred{0}{D_\delta}} \]
  At the real bit the challenge ciphertext decrypts to Bob's input and $\delta$
  is right always; at the zero bit it decrypts to zero and $\delta$ is right
  only on the fiber the leaked output leaves, of mass at most $1/m$. An
  assumption admitting every adversary is therefore forced to assume
  $2\varepsilon_{\mathrm{cpa}} \ge 1 - 1/m$, so the premises of
  Theorem~\ref{thm:alice_guess_admissible} are what make that bound say
  anything.
  Proved in \texttt{dsdp\_alice\_main.v} as
  \texttt{decrypt\_bob\_epsilon\_ge}.
\end{theorem}

\begin{corollary}[A small-epsilon assumption must reject the decryptor]
  \label{cor:decrypt_not_admissible}
  \rocq{infotheo.dumas2017dual.dsdp.hopping.dsdp_alice_main.decrypt_reduction_admissibleF}
  \rocqok
  \uses{thm:decrypt_lower_bound,def:indcpa_assumption}
  With $u_3$ a unit of the plaintext ring, if $\mathcal{A}$ assumes an advantage
  below $1 - 1/m$, its classifier answers false on the reduction adversary the
  decrypting predictor induces at Bob's key.
  \[ \varepsilon_{\mathrm{cpa}} < 1 - \tfrac{1}{m} \implies
     \mathrm{adm}\big(\hopred{0}{D_\delta}\big) = \mathsf{false} \]
  Class membership is therefore not free bookkeeping. On this adversary the
  assumption's own bound decides membership, and decides against it.
  Proved in \texttt{dsdp\_alice\_main.v} as
  \texttt{decrypt\_reduction\_admissibleF}.
\end{corollary}

\begin{corollary}[The premise-free bound is false]
  \label{cor:admissible_premise_free}
  \rocq{infotheo.dumas2017dual.dsdp.hopping.dsdp_alice_main.decrypt_guess_V2_premise_free_lt}
  \rocqok
  \uses{thm:alice_guess_admissible,def:trace_decryptor,def:indcpa_assumption}
  For every assumption whose promised advantage satisfies
  $2\varepsilon_{\mathrm{cpa}} < 1 - 1/m$, the right-hand side of
  Theorem~\ref{thm:alice_guess_admissible} sits strictly below what the
  decrypting predictor achieves.
  \[ \tfrac{1}{m} + 2\,\varepsilon_{\mathrm{cpa}}
       < \Pr[\,\delta(\mathrm{AliceTrace}) = V_2\,] = 1 \]
  Reading that bound without its two membership premises therefore leaves a
  false statement, and for every such assumption rather than a contrived one.
  Corollary~\ref{cor:decrypt_not_admissible} is the complementary half: the
  missing premises cannot be supplied for this predictor.
  Proved in \texttt{dsdp\_alice\_main.v} as
  \texttt{decrypt\_guess\_V2\_premise\_free\_lt}.
\end{corollary}

\begin{definition}[One DSDP instance]
  \label{def:dsdp_instance}
  \rocq{infotheo.dumas2017dual.dsdp.hopping.dsdp_instance.dsdp_instance}
  \rocqok
  \uses{def:alice_trace}
  An instance is an IND-CPA scheme $\mathcal{S}$, the four weights with $u_3$ a
  unit of the plaintext ring, the three private keys, and the two hop coins:
  exactly the parameters every corrupted-Alice trace bound is stated at, so
  every such bound applies at one instance unchanged. The scheme enters as one
  component, so an instance and the assumption made about it name the same
  scheme.
  \[ \mathcal{I} = \big(\mathcal{S},\
       v_1, u_1, u_2, u_3,\ dk_A, dk_B, dk_C,\ w_B, w_C\big) \]
  The real field stays outside the record. One sequence of instances lives over
  one $R$, which only the assumption record and the probabilities mention.
\end{definition}

\begin{definition}[A sequence of DSDP instances]
  \label{def:dsdp_instance_sequence}
  \rocq{infotheo.dumas2017dual.dsdp.hopping.dsdp_instance.dsdp_instance_sequence}
  \rocqok
  \uses{def:dsdp_instance,def:indcpa_assumption}
  A sequence is a map $k \mapsto \mathcal{I}_k$ from the security parameter to
  instances together with an assumption $\mathcal{A}_k$ made at each $k$.
  \[ \mathcal{Q} = \Big(k \mapsto \mathcal{I}_k,\
       k \mapsto \mathcal{A}_k\Big) \]
  The record fixes no relation between consecutive $k$: each instance is
  supplied on its own, and the asymptotic content of a bound along the sequence
  comes from a value of Definition~\ref{def:dsdp_asymptotic} rather than from a
  recurrence between the instances.
\end{definition}

\begin{definition}[The asymptotic content of a sequence]
  \label{def:dsdp_asymptotic}
  \rocq{infotheo.dumas2017dual.dsdp.hopping.dsdp_instance.dsdp_asymptotic}
  \rocqok
  \uses{def:dsdp_instance_sequence,def:negligible_fun}
  An asymptotic value for a sequence $\mathcal{Q}$ is the pair of negligibility
  facts about the two terms of every bound read off along $\mathcal{Q}$.
  \[ \mathcal{N} = \Big(\mathrm{negl}\big(k \mapsto \tfrac{1}{m_k}\big),\
       \mathrm{negl}\big(k \mapsto \varepsilon_{\mathrm{cpa},k}\big)\Big) \]
  The record keeps the two terms apart. The first is unconditional. The inverse
  plaintext cardinality $1/m_k$ is the guessing residue the leaked output
  concedes, it is measured in the plaintext space alone, and it holds against an
  adversary of any running time. The second is conditional on the assumption. It
  is the advantage each $\mathcal{A}_k$ assumes, and it is the only place a
  computational hypothesis enters. The two facts stand apart from $\mathcal{Q}$
  because a statement made at one security parameter needs the instance and the
  assumption there and neither of them, so only the statements that let the
  parameter grow carry an asymptotic value.
\end{definition}

\begin{theorem}[The trace guessing sequence is negligible]
  \label{thm:alice_guess_negligible}
  \rocq{infotheo.dumas2017dual.dsdp.hopping.dsdp_alice_main.alice_trace_guess_V2_negligible,
       infotheo.dumas2017dual.dsdp.hopping.dsdp_alice_main.f_guess_V2,
       infotheo.dumas2017dual.dsdp.hopping.dsdp_alice_main.f_guess_V2_advantageE,
       infotheo.dumas2017dual.dsdp.hopping.dsdp_alice_main.alice_claims_admissible_at,
       infotheo.dumas2017dual.dsdp.hopping.dsdp_alice_main.alice_label_negligible_at}
  \rocqok
  \uses{def:negligible_fun,def:dsdp_instance,def:dsdp_instance_sequence,def:dsdp_asymptotic,thm:alice_guess_admissible,def:alice_trace_chain_admissible,def:negligible_claims,thm:advantage_negligible}
  Along a sequence $\mathcal{Q}$ of DSDP instances, an asymptotic value
  $\mathcal{N}$ for $\mathcal{Q}$, and a sequence of predictors
  $g_k$, if both reduction adversaries of $g_k$ are admissible for
  $\mathcal{A}_k$ at every $k$, then the guessing probabilities are negligible.
  \[ \big(\forall k,\ \mathrm{adm}_{\mathcal{A}_k}(\hopred{0}{D_{g_k}})\big)
     \wedge
     \big(\forall k,\ \mathrm{adm}_{\mathcal{A}_k}(\hopred{1}{D_{g_k}})\big)
     \implies \mathrm{negl}\big(k \mapsto
       \Pr[\,g_k(\mathrm{AliceTrace}_k) = V_2\,]\big) \]
  It is Theorem~\ref{thm:advantage_negligible} read over the program of
  Definition~\ref{def:alice_trace_chain_admissible} at the $k$-th instance. That
  program spends the same three labels at every $k$, and the loss of each of
  them along the sequence is one of the two fields of
  $\mathcal{N}$: both hop labels lose $\mathrm{negl}(k \mapsto
  \varepsilon_{\mathrm{cpa},k})$ and the zero-game label loses
  $\mathrm{negl}(k \mapsto 1/m_k)$. Those two facts are attached to the
  dictionary once, as the negligibleClaims of
  Definition~\ref{def:negligible_claims}, and the theorem names no label.
  The two admissibility facts remain premises, because they restrict the
  reduction adversaries a predictor induces and so speak about the adversary
  rather than about the sequence. That is what separates this statement from
  the decrypting predictor $\delta$, whose guessing probability is $1$ at every
  $k$, and Corollary~\ref{cor:decrypt_sequence_not_admissible} shows that the
  same two fields of $\mathcal{N}$ eventually reject $\delta$'s reduction
  adversary.
\end{theorem}

\begin{corollary}[The sequence eventually rejects the decrypting predictor]
  \label{cor:decrypt_sequence_not_admissible}
  \rocq{infotheo.dumas2017dual.dsdp.hopping.dsdp_alice_main.decrypt_reduction_admissible_eventuallyF}
  \rocqok
  \uses{def:dsdp_instance_sequence,def:dsdp_asymptotic,def:negligible_fun,cor:decrypt_not_admissible}
  Along a sequence $\mathcal{Q}$ with an asymptotic value $\mathcal{N}$ there is
  a threshold past which the classifier
  at each $k$ answers false on the reduction adversary $\delta$ induces at Bob's
  key.
  \[ \exists K,\ \forall k > K,\
     \mathrm{adm}_{\mathcal{A}_k}\big(\hopred{0}{D_\delta}\big)
       = \mathsf{false} \]
  The statement carries no hypothesis beyond $\mathcal{N}$, since its two
  fields supply what the argument needs. At $c = 1$ both fields
  put their sequences below $1/k$, and from $k \ge 2$ that leaves
  $\varepsilon_{\mathrm{cpa},k} < 1 - 1/m_k$, which is the premise of
  Corollary~\ref{cor:decrypt_not_admissible}. The parallel-track reading of
  Theorem~\ref{thm:alice_guess_negligible} is therefore excluded by the very
  sequence that theorem is stated along.
\end{corollary}

\begin{theorem}[The trace simulation distance sequence is negligible]
  \label{thm:alice_trace_sim_advantage_negligible}
  \rocq{infotheo.dumas2017dual.dsdp.hopping.dsdp_alice_main.alice_trace_sim_advantage_negligible,
       infotheo.dumas2017dual.dsdp.hopping.dsdp_alice_main.alice_sim_claims_at,
       infotheo.dumas2017dual.dsdp.hopping.dsdp_alice_main.alice_sim_label_negligible_at}
  \rocqok
  \uses{def:alice_trace_sim_chain_admissible,def:dsdp_asymptotic,def:negligible_claims,thm:advantage_negligible}
  Along a sequence $\mathcal{Q}$ of DSDP instances with an asymptotic value
  $\mathcal{N}$, and a family $D_k$ of tests
  of Alice's executed trace, if both reduction adversaries of $D_k$ are
  admissible for $\mathcal{A}_k$ at every $k$, then the distances between that
  trace and the simulation are negligible.
  \[ \big(\forall k,\ \mathrm{adm}_{\mathcal{A}_k}(\hopred{0}{D_k})\big)
     \wedge
     \big(\forall k,\ \mathrm{adm}_{\mathcal{A}_k}(\hopred{1}{D_k})\big)
     \implies \mathrm{negl}\big(k \mapsto \mathrm{adv}^\tau_k(D_k)\big) \]
  This is computational indistinguishability of Alice's view from the
  simulation, the form a simulation-based secrecy claim takes once the
  parameter is free to grow. It is Theorem~\ref{thm:advantage_negligible} read
  over the program of Definition~\ref{def:alice_trace_sim_chain_admissible} at
  the $k$-th instance, the second statement that program carries. That program
  spends the two hop labels alone, so the whole distance is
  $\mathrm{negl}(k \mapsto \varepsilon_{\mathrm{cpa},k})$ and no
  plaintext-size term enters, which is what separates it from
  Theorem~\ref{thm:alice_guess_negligible}. The dictionary is registered a
  second time, at the family of tests rather than at the test a predictor
  induces, because a statement made at the predictor's own test would be
  weaker than a statement quantified over every test.
\end{theorem}

\begin{definition}[The idealized witnesses]
  \label{def:idealized_setting}
  \rocq{infotheo.dumas2017dual.dsdp.counting.dsdp_random_inputs.uniform_inputs,
       infotheo.computational_security.idealized_indcpa_scheme.idealized_indcpa_scheme,
       infotheo.dumas2017dual.dsdp.hopping.dsdp_alice_instances.idealized_instance,
       infotheo.dumas2017dual.dsdp.hopping.dsdp_alice_instances.idealized_instance_sequence}
  \rocqok
  \uses{def:cipher_constant_assumption,def:negligible_fun}
  Both axes are inhabited at every modulus. The counting witness is the
  uniform law on the five coordinates of $\mathbb{Z}_{pq}$ carrying the three
  plaintext inputs and the two masks, with Alice's three weights held at the
  constants $0$, $0$, $1$ and the three keys constant as well. A weight drawn
  uniformly takes the value $0$ at some sample, and Charlie's weight has to be
  invertible at every sample, which is why the weights stay outside the sample
  space. The hopping witness is \texttt{idealized\_instance\_sequence}, the
  idealized scheme at the composite modulus $p_k q_k$ with $p_k$ a prime above
  $(k+2)^{k+2}$ and $q_k$ a prime above $p_k$, under the cipher-constant
  assumption.
  \[ \varepsilon_{\mathrm{cpa},k} = 0, \qquad
     \#\lvert\mathrm{plain}(\mathcal{S}_k)\rvert = p_k q_k \]
  Every advantage that assumption assumes is zero, so each hopping bound
  reduces at this witness to its unconditional $1/m_k$ term, and each counting
  equality is exact at $\log m_k$. The two factors are taken above
  $(k+2)^{k+2}$ because the modulus growth negligibility asks for has to be
  met by an actual sequence of moduli.
\end{definition}

\begin{definition}[The Paillier instance sequence]
  \label{def:paillier_setting}
  \rocq{infotheo.dumas2017dual.dsdp.hopping.dsdp_alice_instances.paillier_instance,
       infotheo.dumas2017dual.dsdp.hopping.dsdp_alice_instances.paillier_instance_sequence}
  \rocqok
  \uses{def:dsdp_instance_sequence,def:dcr_assumption,def:paillier_indcpa_assumption}
  The sequence whose $k$-th instance is the Paillier scheme at the modulus
  $p_k q_k$ carrying Alice's plaintext input, her three query weights with
  $u_3$ a unit of the plaintext ring, the three private keys and the two
  encryption coins, and whose assumption at $k$ is the IND-CPA record
  Definition~\ref{def:paillier_indcpa_assumption} derives from one decisional
  composite residuosity record $\mathcal{B}_k$.
  \[ \#\lvert\mathrm{plain}(\mathcal{S}_k)\rvert = p_k q_k, \qquad
     \varepsilon_{\mathrm{cpa},k} = 2\,\varepsilon_{\mathrm{dcr},k} \]
  Nothing number-theoretic beyond $1 < p_k, q_k$ enters the sequence itself.
  Primality and distinctness of the two factors are what a Paillier modulus
  is, what decisional composite residuosity starts from, and what the counting
  axis spends on its fiber count, and they are premises of the statements read
  along this sequence rather than fields of it. The three private keys are
  assumed values, no Paillier key being constructed anywhere. What is assumed
  on the computational side is the sequence $\mathcal{B}_k$, so every
  $\varepsilon_{\mathrm{cpa}}$ read off along this sequence is
  $2\varepsilon_{\mathrm{dcr},k}$. The two negligibility facts stay outside and
  enter with the asymptotic value read at it.
\end{definition}

\begin{corollary}[The Paillier readings]
  \label{cor:paillier_security}
  \rocq{infotheo.dumas2017dual.dsdp.hopping.dsdp_alice_instances.paillier_trace_guess_V2_negligible,
       infotheo.dumas2017dual.dsdp.hopping.dsdp_alice_instances.paillier_assumption_at_dcrE,
       infotheo.dumas2017dual.dsdp.hopping.dsdp_alice_instances.paillier_epsilon_at_dcrE,
       infotheo.dumas2017dual.dsdp.hopping.dsdp_alice_instances.paillier_trace_guess_V2_admissible_le,
       infotheo.dumas2017dual.dsdp.hopping.dsdp_alice_instances.paillier_trace_guess_V2_admissible_pq_le}
  \rocqok
  \uses{def:paillier_setting,cor:alice_guess_paillier,thm:alice_guess_negligible,def:dsdp_asymptotic,def:dcr_assumption,def:paillier_indcpa_assumption}
  Along the Paillier sequence a predictor family admissible at every $k$
  matches Bob's input with negligible probability, and at one fixed instance
  the class-conditional trace bound is written in residuosity epsilons, at
  the two spellings of its unconditional term.
  \[ \mathrm{negl}\big(k \mapsto
       \Pr[\,g_k(\mathrm{AliceTrace}_k) = V_2\,]\big), \qquad
     \Pr[\,g(\mathrm{AliceTrace}) = V_2\,]
       \le \tfrac{1}{p q} + 4\,\varepsilon_{\mathrm{dcr}} \]
  The unconditional term is discharged at the plaintext cardinality of the
  Paillier packaging, read as $p q$, with no number theory beyond $1 < p, q$
  entering. The computational term is decisional composite residuosity:
  \texttt{paillier\_assumption\_at\_dcrE} and
  \texttt{paillier\_epsilon\_at\_dcrE} identify the record every
  assumption-conditional bound along the sequence is read at with the one
  derived from $\mathcal{B}_k$, and its advantage with
  $2\varepsilon_{\mathrm{dcr},k}$, both by conversion rather than by a rewrite
  a later statement could route around. The two written-out bounds carry that
  identification into the statement, the factor four being two keys times two
  residuosity calls per key. The negligible statement leaves no
  plaintext-size term visible: that residue enters it as the size half of the
  asymptotic value, one of the two negligible families
  Theorem~\ref{thm:alice_guess_negligible} consumes.
\end{corollary}

\begin{definition}[The Benaloh instance sequence]
  \label{def:benaloh_setting}
  \rocq{infotheo.dumas2017dual.dsdp.hopping.dsdp_alice_instances.benaloh_instance,
       infotheo.dumas2017dual.dsdp.hopping.dsdp_alice_instances.benaloh_instance_sequence}
  \rocqok
  \uses{def:dsdp_instance_sequence,def:benaloh_residuosity_assumption,def:benaloh_indcpa_assumption}
  The sequence whose $k$-th instance is the Benaloh scheme at ciphertext
  modulus $n_k$ and block size $r_k$, carrying the same weights, keys and
  coins as the Paillier sequence, and whose assumption at $k$ is the IND-CPA
  record Definition~\ref{def:benaloh_indcpa_assumption} derives from one
  higher residuosity record $\mathcal{B}_k$ at modulus $n_k$ and exponent
  $r_k$.
  \[ \#\lvert\mathrm{plain}(\mathcal{S}_k)\rvert = r_k, \qquad
     \varepsilon_{\mathrm{cpa},k} = 2\,\varepsilon_{\mathrm{res},k} \]
  Every count is read at the block size $r_k$, which is the plaintext space,
  and not at $n_k$, which sizes the ciphertext space. The sequence asks only
  $1 < r_k$, as the Benaloh scheme, its correctness and its game do. Writing
  $r_k$ as a product of two distinct primes is what the counting axis needs
  for its CRT fiber count, and it enters as a premise of the readings below
  rather than as a field here. Keys are assumed as at Paillier, and the two
  negligibility facts stay outside and enter with the asymptotic value.
\end{definition}

\begin{corollary}[The Benaloh readings]
  \label{cor:benaloh_security}
  \rocq{infotheo.dumas2017dual.dsdp.hopping.dsdp_alice_instances.benaloh_trace_guess_V2_negligible,
       infotheo.dumas2017dual.dsdp.hopping.dsdp_alice_instances.benaloh_assumption_at_residuosityE,
       infotheo.dumas2017dual.dsdp.hopping.dsdp_alice_instances.benaloh_epsilon_at_residuosityE,
       infotheo.dumas2017dual.dsdp.hopping.dsdp_alice_instances.benaloh_trace_guess_V2_admissible_le,
       infotheo.dumas2017dual.dsdp.hopping.dsdp_alice_instances.benaloh_trace_guess_V2_admissible_pq_le}
  \rocqok
  \uses{def:benaloh_setting,cor:alice_guess_benaloh,thm:alice_guess_negligible,def:dsdp_asymptotic,def:benaloh_residuosity_assumption,def:benaloh_indcpa_assumption}
  Along the Benaloh sequence a predictor family admissible at every $k$
  matches Bob's input with negligible probability, and at one fixed instance
  the class-conditional trace bound is written in residuosity epsilons, at
  the two spellings of its unconditional term.
  \[ \mathrm{negl}\big(k \mapsto
       \Pr[\,g_k(\mathrm{AliceTrace}_k) = V_2\,]\big), \qquad
     \Pr[\,g(\mathrm{AliceTrace}) = V_2\,]
       \le \tfrac{1}{r} + 4\,\varepsilon_{\mathrm{res}} \]
  The unconditional term is discharged at the block size $r$, and the second
  spelling reads that block size as $p q$, which is what makes it the same
  number the counting axis measures. The computational term is $r$-th
  residuosity at the ciphertext modulus $n_k$:
  \texttt{benaloh\_assumption\_at\_residuosityE} and
  \texttt{benaloh\_epsilon\_at\_residuosityE} identify the record every
  assumption-conditional bound along the sequence is read at with the one
  derived from $\mathcal{B}_k$, and its advantage with
  $2\varepsilon_{\mathrm{res},k}$, both by conversion. The factor four is two
  keys times two residuosity calls per key, and the unconditional summand
  stays at the block size.
\end{corollary}

\begin{corollary}[The bound is negligible at the idealized witness]
  \label{cor:idealized_security}
  \rocq{infotheo.dumas2017dual.dsdp.hopping.dsdp_alice_instances.alice_trace_guess_V2_idealized_negligible}
  \rocqok
  \uses{def:idealized_setting,thm:alice_guess_negligible,def:cipher_constant_assumption}
  Along the idealized sequence the constant predictor matches Bob's input with
  negligible probability, so the hypotheses of
  Theorem~\ref{thm:alice_guess_negligible} hold together at least once and the
  statement is not empty for want of a sequence and a predictor to read it at.
  \[ \mathrm{negl}\big(k \mapsto
       \Pr[\,0 = V_2\,]\big) \]
  Every advantage the cipher-constant assumption assumes is zero, so the whole
  bound is its unconditional $1/m_k$ term, which is negligible because the
  moduli grow faster than every polynomial. The constant predictor is what the
  cipher-constant class admits: that class keeps the reduction adversaries
  whose verdict ignores the challenge ciphertext, and the decrypting predictor
  is outside it.
\end{corollary}

\begin{theorem}[The trace determines Bob's input]
  \label{thm:centropy_trace_eq0}
  \rocq{infotheo.dumas2017dual.dsdp.hopping.dsdp_alice_trace_link.centropy_V2_trace_eq0}
  \rocqok
  \uses{def:trace_decryptor,thm:centropy_trace_tuple}
  Conditioning on Alice's executed trace leaves no uncertainty about Bob's
  input.
  \[ \HH(V_2 \mid \mathrm{AliceTrace}) = 0 \]
  With Theorem~\ref{thm:centropy_trace_tuple} this transports to her hop-$0$
  tuple and to her full view, so the Shannon restatement of the headline is
  degenerate at the real observation. What carries content there are the
  guessing bounds, which quantify over predictors holding the public keys
  alone; conditional entropy quantifies over nothing and so charges $\delta$ as
  well.
  Proved in \texttt{dsdp\_alice\_trace\_link.v} as
  \texttt{centropy\_V2\_trace\_eq0}.
\end{theorem}

\begin{theorem}[Where the Shannon statement is not degenerate]
  \label{thm:centropy_endpoints}
  \rocq{infotheo.dumas2017dual.dsdp.hopping.dsdp_alice_hop_secrecy.centropy_V2_Sout_logm,
       infotheo.dumas2017dual.dsdp.hopping.dsdp_alice_hop_secrecy.centropy_V2_all_zero_logm}
  \rocqok
  \uses{thm:dsdp_centropy_uniform,def:alice_hop_ladder}
  With $u_3$ a unit of the plaintext ring, conditioning on the leaked output
  alone and conditioning on Alice's all-zero tuple each leave $\log m$ bits of
  uncertainty about Bob's input, which is what it carries unconditioned.
  \[ \HH(V_2 \mid S) = \HH(V_2 \mid \hoptuple{2}) = \log m \]
  At the all-zero endpoint both ciphertext slots encrypt zero, so the leaked
  output is the only channel from $V_2$ into the tuple. This and the $1/m$ term
  of Lemma~\ref{lem:guess_all_zero} descend from one fact, the conditional
  uniformity of $V_2$ given the output; neither is derived from the other.
  Proved in \texttt{dsdp\_alice\_hop\_secrecy.v} as
  \texttt{centropy\_V2\_Sout\_logm} and
  \texttt{centropy\_V2\_all\_zero\_logm}, over the data of one DSDP scheme.
\end{theorem}

\chapter{Tightness: the degenerate query}

The fiber bound of the first chapter needs Alice's third weight $U_3$ to act
invertibly. That hypothesis is tight, not cosmetic: it excludes exactly the query
$U_2 = 1, U_3 = 0$, which is a legal semi-honest input choice, and under it Alice
reads a relay's input outright. These theorems are the boundary of the
corrupted-Alice case, not a separate malicious threat model.

\begin{theorem}[A degenerate query leaks Bob's input]
  \label{thm:US_e1_centropy_V2_eq0}
  \rocq{infotheo.dumas2017dual.dsdp.counting.dsdp_malicious_dotp.US_e1_centropy_V2_eq0}
  \rocqok
  \uses{def:party_views,thm:dsdp_centropy_uniform}
  Fixing her query to the first coordinate ($U_2 = 1$, $U_3 = 0$), Alice reads
  Bob's private input $V_2$ off her view: the protocol output lies in her view and
  equals $V_2$, so its conditional entropy collapses to zero. This is exactly the
  case the fiber bound's injectivity hypothesis excludes, where $1/m$ is false.
  \[ U_2 = 1,\ U_3 = 0 \;\Longrightarrow\; \HH(V_2 \mid \mathrm{AliceView}) = 0 \]
  Proved in \texttt{dsdp\_malicious\_dotp.v} as
  \texttt{US\_e1\_centropy\_V2\_eq0}. The query $U_2 = 1$, $U_3 = 0$ stays a
  premise, being the corrupted Alice's choice and not a property of the sample
  space.
\end{theorem}
\begin{proof}
  \uses{thm:US_e1_centropy_VS0_eq0}
  \rocqok
  The 3-party instance of the N-party leakage theorem at $n=1$
  (Theorem~\ref{thm:US_e1_centropy_VS0_eq0}), with the query the first
  coordinate vector.
\end{proof}

\begin{theorem}[N-party degenerate-query leakage]
  \label{thm:US_e1_centropy_VS0_eq0}
  \rocq{infotheo.dumas2017dual.dsdp.counting.dsdp_malicious_dotp.US_e1_centropy_VS0_eq0}
  \rocqok
  With her query fixed to the first coordinate vector, and the protocol output a
  function of her view, relay party one's input becomes a deterministic image of
  Alice's view, so its conditional entropy is zero. The query is a legal input, so
  this is a semi-honest corruption, not a protocol deviation.
  \[ \mathrm{query} = e_1 \;\Longrightarrow\; \HH(V_{S,0} \mid \mathrm{View}) = 0 \]
\end{theorem}

\part{Corrupted relay}

A corrupted relay forwards masked ciphertexts and never decrypts a peer's input in
the clear. Each relay's whole view is a deterministic image of quantities
independent of the input in question, so conditioning on the view leaves that
input uniform: full $\log m$ bits of uncertainty, hence strictly positive secrecy.
The masking is information-theoretic (a fresh one-time pad), so these bounds carry
no cryptographic assumption.

\chapter{A relay learns nothing about Alice's input}

\begin{theorem}[Relay privacy of $V_1$]
  \label{thm:bob_privacy_V1}
  \rocq{infotheo.dumas2017dual.dsdp.counting.dsdp_relay_secrecy.bob_privacy_V1,
       infotheo.dumas2017dual.dsdp.counting.dsdp_relay_secrecy.charlie_privacy_V1}
  \rocqok
  \uses{def:party_views}
  Assuming Alice's input $V_1$ is a priori uniform and each relay's data is
  independent of $V_1$, the relay's view carries the full $\log m$ bits of
  uncertainty about $V_1$, which is positive. Alice's input enters no message a
  relay sees, so a corrupted relay learns nothing about it.
  \[ \HH(V_1 \mid \mathrm{RelayView}) = \log m > 0 \]
  Proved in \texttt{dsdp\_relay\_secrecy.v} as \texttt{bob\_privacy\_V1}
  and \texttt{charlie\_privacy\_V1}. The uniformity of $V_1$ is a field of the
  inputs record \texttt{dsdp\_random\_inputs}, and the independence of each
  relay's data from $V_1$ is derived there from its each-against-the-rest
  fields.
\end{theorem}
\begin{proof}
  \uses{lem:relay_view_indep_V1}
  \rocqok
  The view is independent of $V_1$ (Lemma~\ref{lem:relay_view_indep_V1}), so
  conditioning does not change the entropy of the a-priori-uniform $V_1$, which is
  $\log m$; positivity is $\log m > \log 1 = 0$.
\end{proof}

\begin{lemma}[The relay view is independent of $V_1$]
  \label{lem:relay_view_indep_V1}
  \rocq{infotheo.dumas2017dual.dsdp.counting.dsdp_relay_secrecy.BobView_indep_V1,
       infotheo.dumas2017dual.dsdp.counting.dsdp_relay_secrecy.CharlieView_indep_V1}
  \rocqok
  Each relay's view is a deterministic image (its key, its input, the forwarded
  ciphertexts) of data assumed independent of $V_1$, hence independent of $V_1$.
  \[ \mathrm{RelayView} \perp V_1 \]
\end{lemma}
\begin{proof}
  \rocqok
  Push the assumed input-independence through the deterministic view map
  (\texttt{inde\_RV\_comp}, an Infotheo library fact).
\end{proof}

\chapter{Bob learns nothing about Charlie's input}

\begin{theorem}[Bob's view hides $V_3$]
  \label{thm:bob_privacy_V3}
  \rocq{infotheo.dumas2017dual.dsdp.counting.dsdp_relay_secrecy.bob_privacy_V3}
  \rocqok
  \uses{def:party_views}
  With $V_3$ and the mask $R_3$ a priori uniform and independent as stated, Bob's
  view carries the full $\log m$ bits of uncertainty about Charlie's input $V_3$.
  The Charlie-ciphertext Bob forwards carries $V_3 U_3 + R_3$, masked by Alice's
  fresh one-time pad $R_3$, so Bob's view is independent of $V_3$.
  \[ \HH(V_3 \mid \mathrm{BobView}) = \log m > 0 \]
  Proved in \texttt{dsdp\_relay\_secrecy.v} as \texttt{bob\_privacy\_V3}.
  The uniformity of $V_3$ and of $R_3$ are fields of the inputs record
  \texttt{dsdp\_random\_inputs}, and the freshness of $R_3$ against Charlie's
  weighted input is derived there from its each-against-the-rest fields.
\end{theorem}
\begin{proof}
  \uses{lem:one_time_pad}
  \rocqok
  The masked term $V_3 U_3 + R_3$ is independent of $V_3$ by the one-time-pad
  lemma (Lemma~\ref{lem:one_time_pad}); assembling this with Bob's clean data
  through the graphoid mixing rule gives $\mathrm{BobView} \perp V_3$, and the
  entropy computation follows as before.
\end{proof}

\chapter{Charlie learns nothing about Bob's input}

\begin{theorem}[Charlie's view hides $V_2$]
  \label{thm:charlie_privacy_V2}
  \rocq{infotheo.dumas2017dual.dsdp.counting.dsdp_relay_secrecy.charlie_privacy_V2}
  \rocqok
  \uses{def:party_views}
  With $V_2$ and the mask $R_2$ a priori uniform and independent as stated,
  Charlie's view carries the full $\log m$ bits of uncertainty about Bob's input
  $V_2$. Charlie's decrypted aggregate $D_3$ carries $V_2 U_2$ masked by Alice's
  fresh pad $R_2$, so the ciphertext Charlie returns to Alice is independent of
  $V_2$.
  \[ \HH(V_2 \mid \mathrm{CharlieView}) = \log m > 0 \]
  Proved in \texttt{dsdp\_relay\_secrecy.v} as
  \texttt{charlie\_privacy\_V2}. The uniformity of $V_2$ and of $R_2$ are
  fields of the inputs record \texttt{dsdp\_random\_inputs}, and the freshness
  of $R_2$ against Bob's weighted input is derived there from its
  each-against-the-rest fields.
\end{theorem}
\begin{proof}
  \uses{lem:one_time_pad}
  \rocqok
  The pad $R_2$ hides $V_2 U_2$; a second pad argument passes $D_3 = V_2 U_2 + R_2
  + \dots$ to independence of $V_2$, and the mixing rule assembles Charlie's full
  view, leaving $V_2$ uniform.
\end{proof}

\begin{lemma}[One-time-pad masking]
  \label{lem:one_time_pad}
  \rocqok
  % library: lemma_3_5' (the du2002 one-time-pad / secret-sharing fact), cited as
  %   a black-box leaf. Used by the two relay-secrecy chapters above.
  A value masked by a fresh uniform pad independent of it is independent of that
  value: $X + R$ with $R$ uniform and $R \perp X$ makes $X + R$ independent of
  $X$. This is the information-theoretic core of the relay-secrecy bounds and
  needs no cryptographic assumption.
  \[ R \text{ uniform},\ R \perp X \;\Longrightarrow\; (X + R) \perp X \]
\end{lemma}
